% =============================================================
% BKR_FFT_Teslja — main.tex
% Бакалаврська робота: Візуальне моделювання ШПФ
% НУ «Львівська політехніка», каф. АСУ, 2025
% =============================================================

\documentclass[14pt,a4paper]{extarticle}

% --- Мова і кодування ---
\usepackage[utf8]{inputenc}
\usepackage[ukrainian,english]{babel}

% --- Шрифт Times New Roman ---
\usepackage{tempora}
\usepackage[T1,T2A]{fontenc}


% --- Поля (вимоги кафедри АСУ) ---
\usepackage[
  left=25mm,
  right=10mm,
  top=20mm,
  bottom=20mm
]{geometry}

% --- Міжрядковий інтервал 1.5 ---
\usepackage{setspace}
\onehalfspacing

% --- Математика ---
\usepackage{amsmath}
\usepackage{amssymb}
\usepackage{amsfonts}

% --- Зображення ---
\usepackage{graphicx}
\usepackage{float}
\graphicspath{{images/}}

% --- Таблиці ---
\usepackage{booktabs}
\usepackage{array}
\usepackage{longtable}
\usepackage{multirow}

% --- Код і лістинги ---
\usepackage{listings}
\usepackage{xcolor}

\lstdefinestyle{jsstyle}{
  language=Java,
  basicstyle=\small\ttfamily,
  keywordstyle=\color{blue}\bfseries,
  commentstyle=\color{gray}\itshape,
  stringstyle=\color{orange},
  numberstyle=\tiny\color{gray},
  numbers=left,
  numbersep=8pt,
  breaklines=true,
  breakatwhitespace=true,
  frame=single,
  rulecolor=\color{gray!30},
  tabsize=2
}

\lstdefinestyle{cstyle}{
  language=C,
  basicstyle=\small\ttfamily,
  keywordstyle=\color{blue}\bfseries,
  commentstyle=\color{gray}\itshape,
  stringstyle=\color{orange},
  numberstyle=\tiny\color{gray},
  numbers=left,
  numbersep=8pt,
  breaklines=true,
  frame=single,
  rulecolor=\color{gray!30},
  tabsize=2
}

% --- Гіперпосилання ---
\usepackage{hyperref}
\hypersetup{
  colorlinks=true,
  linkcolor=black,
  citecolor=black,
  urlcolor=blue,
  pdfauthor={Тесля М.},
  pdftitle={Візуальне моделювання роботи обчислювача ШПФ під керуванням потоку даних}
}

% --- Нумерація рисунків, таблиць, формул у межах розділу ---
\usepackage{chngcntr}
\counterwithin{figure}{section}
\counterwithin{table}{section}
\counterwithin{equation}{section}

% --- Бібліографія ---
\usepackage[square,numbers,sort]{natbib}

% --- Відступи і параграфи ---
\usepackage{indentfirst}
\setlength{\parindent}{1.25cm}

% --- Налаштування заголовків ---
\usepackage{titlesec}
\titleformat{\section}[block]
  {\normalfont\large\bfseries\centering}
  {РОЗДІЛ \thesection.}{0.5em}{}
\titleformat{\subsection}[block]
  {\normalfont\normalsize\bfseries}
  {\thesubsection}{0.5em}{}
\titleformat{\subsubsection}[block]
  {\normalfont\normalsize\bfseries}
  {\thesubsubsection}{0.5em}{}

% --- Зміст ---
\usepackage{tocloft}
\renewcommand{\cftsecleader}{\cftdotfill{\cftdotsep}}
\usepackage{lastpage}

% =============================================================
\begin{document}

% --- Титульний аркуш ---
% (вставляється окремим PDF або вручну друкується на бланку)

% --- Завдання ---
% (вставляється на бланку кафедри)

% --- Анотація ---
% chapters/abstract.tex
\addcontentsline{toc}{section}{АНОТАЦІЯ}
{\centering\bfseries АНОТАЦІЯ\par}
\vspace{1em}

\textbf{Тесля М.} Візуальне моделювання роботи обчислювача ШПФ під
керуванням потоку даних. Бакалаврська кваліфікаційна робота за
спеціальністю 122 «Комп'ютерні науки». Національний університет
«Львівська політехніка», Львів, 2026.

У межах бакалаврської кваліфікаційної роботи розроблено математичну
модель, граф потоку даних та веб-орієнтований програмний симулятор
16-точкового швидкого перетворення Фур'є за алгоритмом radix-4 для
інтерактивної візуалізації потокового виконання обчислень із
використанням технологій Next.js, React~Flow, Zustand і TypeScript.

Робота складається із вступу, чотирьох основних розділів, економічної
частини, висновків, списку використаних джерел та додатків.

У вступі обґрунтовано актуальність теми, визначено мету, задачі,
об'єкт і предмет дослідження, описано методологію та інструментарій,
а також практичну цінність розробки.

У першому розділі проведено аналіз предметної області, огляд
сучасних підходів до обчислення ШПФ, порівняльний аналіз існуючих
рішень (FFTW, CMSIS-DSP, cuFFT, SPIRAL), виконано системний аналіз
із виділенням зацікавлених сторін, цілей та обмежень, розглянуто
етичні й правові аспекти та сформульовано формальну постановку
задачі з вимогами F1--F9, NF1--NF8 і сценаріями використання UC1--UC9.

Другий розділ присвячено проєктуванню системи: обґрунтовано вибір
клієнтського SPA-підходу, розроблено математичну модель тристадійної
$4\times4$ декомпозиції, сформовано граф потоку даних (56~вузлів,
64~дуги), визначено критичний шлях і коефіцієнт паралелізму, описано
структури даних та механізми обміну між компонентами.

Третій розділ описує програмну реалізацію симулятора: технологічний
стек, структуру проєкту, реалізацію ядра \texttt{DataflowRuntime}
з методом \texttt{tick()}, обчислювальних вузлів, алгоритму
\texttt{generateDFT4x4()}, класифікацію винятків, процедуру
розгортання та модульне тестування у середовищі Jest.

У четвертому розділі подано результати експериментальних досліджень:
верифікацію на шести тестових сценаріях (максимальна похибка
$3\times10^{-12}$ при допуску $\varepsilon=10^{-9}$), аналіз
обчислювальної складності ($K_M\approx28{,}4\times$), аналіз
паралелізму dataflow-моделі ($K_p\approx1{,}88\times$),
обговорення, моніторинг експлуатаційних метрик та підтвердження
нефункціональних вимог NF1--NF8.

У висновках узагальнено основні результати виконаної роботи,
визначено практичне значення та перспективні напрями подальших
досліджень.

Дана бакалаврська кваліфікаційна робота містить: \pageref{LastPage}~сторінок,
1~додаток та 25~джерел літератури.

\textbf{Ключові слова:} швидке перетворення Фур'є, radix-4, граф
потоку даних, dataflow computing, візуальне моделювання,
веб-симулятор, Next.js, React~Flow, паралелізм, цифрове оброблення
сигналів.

\vspace{1em}\hrule\vspace{1em}

\addcontentsline{toc}{section}{ABSTRACT}
{\centering\bfseries ABSTRACT\par}
\vspace{1em}

\begin{otherlanguage}{english}
\textbf{Teslia M.} Visual Modeling of a Dataflow-Driven FFT Processor.
Bachelor's Qualification Work for specialty 122 «Computer Science».
Lviv Polytechnic National University, Lviv, 2026.

Within this bachelor's qualification work, a mathematical model, a
data-flow graph, and a web-based software simulator of a 16-point
Fast Fourier Transform (radix-4) have been developed for interactive
visualization of the dataflow-driven computation, using Next.js,
React~Flow, Zustand, and TypeScript technologies.

The work consists of an introduction, four main chapters, an economic
part, conclusions, a list of references, and appendices.

The introduction substantiates the relevance of the topic and defines
the aim, tasks, object, subject of the study, the methodology, and
the practical value of the development.

The first chapter presents an analysis of the problem domain, a
review of modern FFT computation approaches, a comparative analysis
of existing solutions (FFTW, CMSIS-DSP, cuFFT, SPIRAL), a systems
analysis with stakeholders, goals and constraints, a discussion of
ethical and legal aspects, and a formal task statement with F1--F9,
NF1--NF8 requirements and UC1--UC9 use cases.

The second chapter is devoted to the system design: justification of
a client-side SPA approach, a mathematical model of the three-stage
$4\times4$ decomposition, a data-flow graph (56~nodes, 64~arcs),
determination of the critical path and parallelism factor, data
structures, and inter-component communication mechanisms.

The third chapter describes the software implementation of the
simulator: technology stack, project structure, the
\texttt{DataflowRuntime} engine with the \texttt{tick()} method,
computational nodes, the \texttt{generateDFT4x4()} algorithm,
exception handling, deployment procedure, and unit testing using Jest.

The fourth chapter presents the experimental results: verification
across six test scenarios (maximum error $3\times10^{-12}$ against
tolerance $\varepsilon=10^{-9}$), analysis of computational complexity
($K_M\approx28.4\times$), analysis of dataflow-model parallelism
($K_p\approx1.88\times$), discussion, monitoring of operational
metrics, and confirmation of non-functional requirements NF1--NF8.

The conclusions summarize the key results of the work and outline
the practical value and promising directions for further research.

This bachelor's qualification work contains: \pageref{LastPage}~pages,
1~appendix, and 25~references.

\textbf{Keywords:} fast Fourier transform, radix-4, data-flow graph,
dataflow computing, visual modeling, web-based simulator, Next.js,
React~Flow, parallelism, digital signal processing.
\end{otherlanguage}

\newpage

% --- Зміст ---
\tableofcontents
\newpage

% --- Перелік скорочень ---
% chapters/abbreviations.tex
\addcontentsline{toc}{section}{ПЕРЕЛІК УМОВНИХ СКОРОЧЕНЬ, СИМВОЛІВ І ТЕРМІНІВ}
{\centering\bfseries ПЕРЕЛІК УМОВНИХ СКОРОЧЕНЬ, СИМВОЛІВ І ТЕРМІНІВ\par}
\vspace{1em}

\begin{tabular}{p{2.5cm}p{11.5cm}}
  АЦП    & аналого-цифровий перетворювач \\
  БКР    & бакалаврська кваліфікаційна робота \\
  ДПФ    & дискретне перетворення Фур'є \\
  ЕЕГ    & електроенцефалограма \\
  КС     & керуюче слово \\
  МРТ    & магнітно-резонансна томографія \\
  ОП     & оперативна пам'ять \\
  ПЗ     & програмне забезпечення \\
  ПК     & персональний комп'ютер \\
  ЦОС    & цифрове оброблення сигналів \\
  ШПФ    & швидке перетворення Фур'є \\
  \hline
  API    & Application Programming Interface (прикладний програмний інтерфейс) \\
  C4     & Context--Container--Component--Code (нотація опису архітектури) \\
  CI/CD  & Continuous Integration / Continuous Delivery \\
  DAG    & Directed Acyclic Graph (спрямований ациклічний граф) \\
  DFG    & Data Flow Graph (граф потоку даних) \\
  DFT    & Discrete Fourier Transform \\
  DSP    & Digital Signal Processing \\
  EMA    & Exponential Moving Average (експоненційне ковзне середнє) \\
  F1--F9 & функціональні вимоги до системи \\
  FFT    & Fast Fourier Transform \\
  FIFO   & First In, First Out \\
  FPGA   & Field-Programmable Gate Array \\
  GDPR   & General Data Protection Regulation \\
  IEEE   & Institute of Electrical and Electronics Engineers \\
  IEEE 754 & стандарт подання чисел з плаваючою крапкою \\
  LMS    & Learning Management System (система управління навчанням) \\
  LTI    & Learning Tools Interoperability \\
  MIT    & Massachusetts Institute of Technology (ліцензія з відкритим кодом) \\
  NF1--NF8 & нефункціональні вимоги до системи \\
  OFDM   & Orthogonal Frequency-Division Multiplexing \\
  RAM    & Random Access Memory \\
  rAF    & requestAnimationFrame \\
  RTL    & Register-Transfer Level (рівень регістрових передач) \\
  SDF    & Synchronous Data Flow (синхронний потік даних) \\
  SPA    & Single Page Application (односторінковий застосунок) \\
  SSR    & Server-Side Rendering \\
  TF     & Twiddle Factor (поворотний множник) \\
  UC1--UC9 & Use Cases (сценарії використання) \\
  UI     & User Interface (інтерфейс користувача) \\
  UML    & Unified Modeling Language \\
  UX     & User Experience (користувацький досвід) \\
  XSS    & Cross-Site Scripting \\
\end{tabular}

\newpage

% --- Вступ ---
% chapters/intro.tex
\addcontentsline{toc}{section}{ВСТУП}
{\centering\bfseries ВСТУП\par}
\vspace{1em}

\textbf{Актуальність теми.}
Швидке перетворення Фур'є (ШПФ, англ.\ \textit{Fast Fourier Transform}) належить
до фундаментальних алгоритмів цифрового оброблення сигналів (ЦОС) і щоденно
виконується мільярди разів у мобільному зв'язку 4G/5G (OFDM-модуляція), системах
активного шумозаглушення, медичній діагностиці (МРТ, ЕЕГ-спектрометрія) та
обробленні аудіо- і відеопотоків. За оцінками IEEE Signal Processing Society
(2024), ШПФ лишається найчастіше застосовуваним числовим алгоритмом у вбудованих
системах реального часу. Алгоритм radix-4 для $N = 16$ через тристадійну
$4 \times 4$ декомпозицію потребує лише 9 нетривіальних комплексних множень
проти 256 у прямому ДПФ, що дає прискорення у 28~разів.

Попри наявність зрілих виробничих бібліотек (FFTW, ARM CMSIS-DSP, NVIDIA cuFFT,
Intel IPP), жодна з них не надає навчально-дослідницького засобу для
\emph{безпосереднього спостереження} потокової природи алгоритму: динаміки
токенів між вузлами, паралельної активації незалежних обчислювальних блоків та
наскрізної затримки у графі обчислень. Існуючі засоби відображають лише
числовий результат, приховуючи мікроархітектурний хід обчислення. Водночас
парадигма обчислень під керуванням потоку даних (dataflow computing) є природною
для ШПФ: граф метеликів безпосередньо відображається на граф потоку даних
(DFG), а чотири незалежні вузли 4-точкового ДПФ кожного ярусу активуються
паралельно без явної синхронізації. Це зумовлює актуальність розроблення
веб-орієнтованого симулятора, що поєднує достовірне виконання реального
алгоритму radix-4 з інтерактивною DFG-візуалізацією.

\textbf{Мета роботи} --- розроблення математичної моделі, графу потоку даних та
програмного симулятора 16-точкового ШПФ з основою 4 для інтерактивної
візуалізації потокового виконання алгоритму у веб-браузері.

\textbf{Об'єктом дослідження} є процеси обчислення дискретного перетворення
Фур'є під керуванням потоку даних у програмних симуляторах обчислювальних систем.

\textbf{Предметом дослідження} є методи та підходи до математичного моделювання
алгоритму radix-4, побудови графу потоку даних (DFG / SDF) обчислювача та
програмної реалізації інтерактивного симулятора на платформі веб-браузера.

\textbf{Задачі дослідження:}
\begin{enumerate}
  \item проаналізувати сімейство алгоритмів ШПФ, парадигму dataflow computing та
    методи візуального моделювання обчислювальних систем; виконати системний
    аналіз предметної області та обґрунтувати вибір архітектурного підходу;
  \item розробити математичну модель алгоритму radix-4 для $N = 16$ з
    відображенням індексів, трьома стадіями обчислення, матрицею поворотних
    множників та оцінкою обчислювальної складності;
  \item розробити граф потоку даних обчислювача зі специфікацією вузлів, дуг,
    затримок і керуючих слів; визначити критичний шлях та коефіцієнт
    паралелізму DFG;
  \item реалізувати ядро симуляції \texttt{DataflowRuntime} з правилом активації
    токенів та веб-інтерфейс візуалізації на React Flow з анімацією потоку
    токенів і тепловою картою активності;
  \item виконати верифікацію результатів симулятора порівнянням із еталонним ДПФ
    бібліотеки \texttt{fft.js} із точністю $\varepsilon = 10^{-9}$ на шести
    тестових сценаріях;
  \item провести експериментальний аналіз обчислювальної складності та
    паралелізму dataflow-моделі; підтвердити теоретичні оцінки рантаймовими
    метриками симулятора.
\end{enumerate}

\textbf{Методологія та інструментарій.} Математична база дослідження --- теорія
дискретних перетворень Фур'є та алгоритмів Cooley--Tukey, формалізм синхронних
dataflow-графів (SDF) Е.~Лі і Е.~Мессершмітта. Інженерна основа --- архітектурне
моделювання за нотацією C4 та UML-діаграмами класів і послідовностей. Програмна
реалізація використовує стек Next.js~14 / React~18 / TypeScript~5 із бібліотеками
React Flow (\texttt{@xyflow/react}) 12 для рендерингу графу, Zustand~4 для
реактивного управління станом, \texttt{fft.js} як еталон верифікації та Jest~29
для модульного й інтеграційного тестування.

\textbf{Практична цінність роботи.} Розроблений симулятор є
навчально-дослідницьким інструментом для вивчення потокових алгоритмів ЦОС у
курсі «Цифрове оброблення сигналів» та суміжних дисциплінах бакалаврської
програми 122 «Комп'ютерні науки». Застосунок виконує реальний алгоритм radix-4
(не його імітацію), доступний без встановлення додаткового програмного
забезпечення через веб-браузер і підтримує регульовану швидкість симуляції та
покроковий режим для прецизійного дослідження активацій вузлів. Верифікована
точність $\varepsilon < 10^{-9}$ підтверджує придатність для прикладних числових
досліджень. Архітектура DFG (56 вузлів, 64 дуги, параметри \texttt{latency} і
\texttt{delay}) може слугувати вихідною специфікацією для апаратної реалізації
обчислювача на FPGA, де граф потоку даних безпосередньо відображається на
маршрутизацію сигналів між DSP-блоками.

\textbf{Апробація результатів.} Окремих публікацій за темою роботи не подавалось;
результати роботи планується подати на студентську науково-технічну конференцію
НУ «Львівська політехніка» у 2026~році.
\newpage

% --- Розділ 1 ---
% chapters/chapter1.tex
\section{АНАЛІЗ ПРЕДМЕТНОЇ ОБЛАСТІ ТА ПОСТАНОВКА ЗАДАЧІ}

% Розділ 1 згідно нових методичних вказівок (10-15 стор.):
% 1.1 Характеристика предметної області
% 1.2 Огляд сучасних підходів та рішень
% 1.3 Системний аналіз задачі
% 1.4 Етичні, правові та регуляторні аспекти
% 1.5 Постановка задачі та формулювання вимог
% 1.6 Висновки до розділу 1
%
% Математичне та алгоритмічне обґрунтування (ДПФ, radix-4, DFG-модель,
% матричні факторизації) перенесено до розділу 2 «Проєктування».

% ─────────────────────────────────────────────────
\subsection{Характеристика предметної області}

\subsubsection*{Контекст застосування}

Швидке перетворення Фур'є (ШПФ, англ.\ \textit{Fast Fourier Transform}
--- FFT) є одним із найбільш вживаних обчислювальних алгоритмів
цифрової обробки сигналів (ЦОС). Воно лежить в основі цифрової
телекомунікації (OFDM-модуляція у стандартах LTE, 5G, Wi-Fi),
обробки мови й звуку, радіолокації, медичного томографування,
спектрального аналізу та сучасних алгоритмів машинного навчання
з перетворенням вхідних даних у частотний простір~\cite{nussbaumer1982,protsko2025}.
За обсягом щоденно виконуваних операцій ШПФ залишається однією з
найзавантаженіших обчислювальних процедур у світі.

Пряме обчислення дискретного перетворення Фур'є (ДПФ) має
квадратичну складність $O(N^2)$. Навіть для скромного розміру
$N = 1024$ це означає понад мільйон комплексних множень,
що неприйнятно для задач реального часу. Клас швидких алгоритмів,
започаткований Дж.~Кулі та Дж.~Тюкі у 1965~р.~\cite{cooley1965},
знижує складність до $O(N \log_2 N)$ завдяки рекурсивній
факторизації базисної матриці~\cite{nussbaumer1982}. Для довжин,
що є степенями чотирьох ($N = 4^t$), найефективнішим з точки зору
кількості нетривіальних множень є алгоритм з основою~4 (radix-4),
який шляхом двовимірної декомпозиції зводить одновимірне ДПФ
довжиною $N = 16$ до тристадійної схеми $4 \times 4$: чотири
паралельних ДПФ-4, множення на матрицю поворотних множників
$W_{16}^{rk}$ та знову чотири паралельних ДПФ-4~\cite{burrus1985dft}.

Паралельна природа алгоритму природно вкладається у парадигму
обчислень під керуванням потоку даних (англ.\ \textit{dataflow
computing}), теоретичні засади якої були закладені Дж.~Деннісом
(Jack B.\ Dennis) у MIT на початку 1970-х~рр.~\cite{dennis1974,dennis1975}.
У dataflow-моделі відсутній лічильник команд фон-нейманівської
архітектури: вузол графа активується тоді і лише тоді, коли на
всіх його вхідних каналах є токени, після чого самопланування
системи (англ.\ \textit{self-scheduling}) забезпечує максимальний
паралелізм без явних механізмів синхронізації~\cite{veen1986}.
Граф алгоритму ШПФ, що формально є орієнтованим ациклічним графом
залежностей за даними, відображається на граф потоку даних
(англ.\ \textit{Data Flow Graph} --- DFG) без перетворення моделі.

\subsubsection*{Опис проблеми та її актуальність}

Незважаючи на фундаментальну роль ШПФ та широкий спектр його
виробничих реалізацій, у навчальній та дослідницькій практиці
існує системна прогалина між трьома рівнями опису алгоритму:

\begin{itemize}
  \item \textit{математична модель} --- матричні факторизації,
    формули метеликів і поворотних множників, що подаються у
    підручниках із ЦОС~\cite{nussbaumer1982,burrus1985dft};
  \item \textit{формальна потокова модель} --- DFG і Synchronous
    Dataflow (SDF) Лі та Месершмітта, що описують вузли, дуги,
    токени та правило активації~\cite{lee1994,veen1986};
  \item \textit{виробнича реалізація} --- бібліотеки FFTW,
    ARM~CMSIS-DSP, NVIDIA~CuFFT, SPIRAL, Intel~IPP, що оптимізовані
    виключно під максимальну продуктивність на конкретних
    платформах~\cite{frigo2005,cmsis_dsp,du2012,puschel2004,intel_ipp}.
\end{itemize}

Ці рівні існують відокремлено один від одного. Математичні формули
не мають виконуваного аналога з можливістю поточного спостереження;
формальна SDF-модель рідко реалізується як самостійний артефакт;
а виробничі бібліотеки не надають жодних засобів для інтерактивного
дослідження внутрішньої потокової природи алгоритму: процесу
породження та поглинання токенів, динаміки завантаження вузлів,
часу наскрізної затримки та паралельних шляхів у графі. Студент,
викладач або інженер-початківець, що вивчає архітектуру потокового
обчислювача ШПФ, не має інструменту, який одночасно є математично
достовірним (виконує реальний алгоритм), формально обґрунтованим
(відповідає SDF-моделі) і візуально інтерактивним (показує рух
токенів у реальному часі).

Актуальність розробки візуального симулятора потокового обчислювача
ШПФ зумовлена трьома чинниками: (1)~зростанням попиту на навчальні
засоби з інтерактивною візуалізацією обчислювальних процесів у
підготовці ІТ-фахівців; (2)~розширенням інтересу до dataflow- й
потокових архітектур у сучасних спеціалізованих прискорювачах
(зокрема, тензорних процесорах і FPGA-рішеннях для ЦОС);
(3)~відсутністю вітчизняних веб-орієнтованих інструментів із
відкритим кодом для цього класу задач.


% ─────────────────────────────────────────────────
\subsection{Огляд сучасних підходів та рішень}

Аналіз літературних джерел і наявних програмних рішень за останні
10~років показує, що простір сучасних засобів обчислення ШПФ
поділяється на три сегменти: (а)~програмні бібліотеки загального
призначення; (б)~апаратні та GPU-прискорювачі; (в)~системи
автоматичної генерації оптимізованого коду. Паралельно існує
окрема галузь методів візуального моделювання, які застосовуються
для формалізованого опису обчислювачів, однак не реалізуються як
виконувані симулятори. Нижче розглядаються обидва ці напрямки та
проводиться їх порівняльний аналіз.

\subsubsection*{Програмні бібліотеки обчислення ШПФ}

\textbf{Бібліотека FFTW (Fastest Fourier Transform in the West)}
розроблена М.~Фріго і С.~Джонсоном у MIT у 1997~р.~\cite{frigo1998}
і у~1999~р. відзначена премією Вілкінсона. Архітектура FFTW
ґрунтується на трьох компонентах~\cite{frigo2005}: \emph{коделети}
--- компактні блоки прямолінійного C-коду для ДПФ невеликих розмірів,
автоматично згенеровані компілятором \texttt{genfft}; \emph{планувальник}
--- runtime-компонент, що методом динамічного програмування добирає
оптимальну рекурсивну факторизацію під конкретну апаратну платформу;
\emph{екзекутор} --- виконує складений план через явну рекурсію
«розділяй і володарюй». Рекурсивна декомпозиція є кеш-незалежною
(cache-oblivious): щойно підзадача поміщається у L1/L2-кеш, вона
виконується без промахів кешу, що і забезпечує вищу продуктивність
проти ітеративних реалізацій~\cite{frigo2005}.

\textbf{ARM CMSIS-DSP} --- стандартизований пакет ЦОС для
мікроконтролерів ARM Cortex-M~\cite{cmsis_dsp}. Надає жорстко
оптимізовані реалізації radix-4 для розмірів $N = 16, 64, 256,
1024, 4096$ із ручною адаптацією під SIMD-інструкції NEON і
DSP-розширення Cortex-M4/M7. Бібліотека не переноситься на інші
архітектури, не підтримує адаптивного планування та не надає
будь-яких засобів моніторингу стану обчислення.

\textbf{NVIDIA CuFFT} та \textbf{OpenCL FFT} --- адаптації ШПФ
для графічних процесорів. Стандарт OpenCL забезпечує функціональну
портативність, але не гарантує продуктивності: ядро, оптимізоване
під архітектуру NVIDIA~Fermi, демонструє низьку ефективність на
GPU ATI~Radeon через відмінності в ієрархії пам'яті та механізмах
планування потоків (warps/wavefronts)~\cite{du2012}. Для подолання
цього застосовується автотюнінг: А.~Нукада та С.~Мацуока~\cite{nukada2009}
показали, що адаптивний вибір відступів у спільній пам'яті та
кількості активних блоків потоків дозволяє перевершити навіть
вендорну бібліотеку NVIDIA~CuFFT.

\textbf{SPIRAL} (Carnegie Mellon University) --- принципово інший
підхід: замість вибору між написаними ядрами, SPIRAL генерує
платформо-залежний код безпосередньо з математичних формул
факторизації, описаних мовою SPL (Signal Processing Language).
Пошук оптимального графа зводиться до задачі пошуку у просторі
еквівалентних матричних розкладів, що може вирішуватися евристично
або методами машинного навчання~\cite{puschel2004}. SPIRAL
генерує код як для CPU (SIMD, AVX), так і для FPGA.

\textbf{Intel IPP (Integrated Performance Primitives)} займає
проміжне місце: жорстко оптимізовані примітиви ШПФ під векторні
інструкції SSE/AVX/AVX-512 процесорів Intel без автоматичного
пошуку факторизації під час виконання~\cite{intel_ipp}.

\subsubsection*{Методи візуального моделювання обчислювальних систем}

Окрему групу джерел утворюють методи візуального моделювання,
які теоретично описують структуру обчислювачів, але не реалізуються
як виконувані симулятори з числовою верифікацією.

\textbf{Загальні нотації моделювання} --- UML~Activity,
IDEF0/IDEF3, DFD Йордана/ДеМарко --- орієнтовані на control-flow
парадигму~\cite{rumbaugh2004,yourdon1979}. Вони не мають формальної
семантики токенів і не забезпечують однозначного відображення
паралелізму за даними, тому для потокового обчислювача ШПФ
породжують надлишкові конструкції \texttt{fork}/\texttt{join}.

\textbf{DFG/SDF} (Data Flow Graph і Synchronous Dataflow) ---
спеціалізовані моделі, розроблені для задач ЦОС~\cite{veen1986,lee1994,mit_dataflow}.
У SDF-моделі кількість токенів, що споживаються та виробляються
кожним вузлом за одну активацію, є фіксованою і відомою на етапі
статичного аналізу. Це дозволяє виконувати статичне планування без
ризику взаємних блокувань, гарантувати детермінованість результатів
і перевіряти консистентність графа алгебраїчно через топологічну
матрицю~\cite{schaumont2019}. Алгоритм ШПФ із фіксованою довжиною
$N = 16$ є природним прикладом SDF-графа.

\textbf{Сигнальні графи (signal flowgraphs) і butterfly-діаграми}
--- С.~С.~Буррус формалізує зв'язок між матричною факторизацією та
графом алгоритму через рівняння $\mathbf{X'} = \mathbf{C} \cdot
\mathbf{D} \cdot \mathbf{A} \cdot \mathbf{X}$, де кожна матриця
відповідає окремому ярусу графа~\cite{burrus1985dft}. Ці діаграми
є статичними і не моделюють динаміку виконання. У вітчизняній
науковій традиції аналогічний підхід застосовувався у роботах
ОКБ~ЕІВТ (М.~Розенблат, 1970-ті~рр.) та у монографії М.~М.~Яцимірського
(1997~р.)~\cite{protsko2025}. Зокрема, у Фізико-механічному
інституті ім.~Г.~В.~Карпенка НАН~України в 1980-х роках було
розроблено векторний спецпроцесор ШПФ на основі саме dataflow-парадигми
з блоком арифметики, що апаратно реалізовував чотири базові
потокові операції~\cite{protsko2025}.

\subsubsection*{Ступінь дослідженості проблеми та вузькі місця}

Узагальнення розглянутих джерел дозволяє зафіксувати такі
прогалини. По-перше, виробничі бібліотеки (FFTW, CMSIS-DSP,
CuFFT, SPIRAL, Intel~IPP) оптимізовані під продуктивність і не
надають жодних інтроспективних засобів: FFTW розкриває план
факторизації лише у текстовому форматі, жодна з бібліотек не
генерує анімований DFG і не показує рух токенів. По-друге, теоретичні
моделі DFG/SDF формально визначені, проте рідко супроводжуються
веб-доступним симулятором із числовою верифікацією. По-третє,
статичні засоби (draw.io, Visio, Simulink) або не відтворюють
динаміки, або потребують комерційних ліцензій, що обмежує
відтворюваність у навчальному та дослідницькому контекстах.

Таким чином, більшість авторів вирішує задачу обчислення ШПФ або
суто продуктивнісним підходом (FFTW, CuFFT, SPIRAL), або суто
теоретичним (DFG/SDF-моделі, сигнальні графи), однак підхід, який
поєднує формальну SDF-модель, виконуваний алгоритм із числовою
верифікацією та інтерактивну веб-візуалізацію, у літературі не
представлений.

\subsubsection*{Порівняння існуючих рішень}

Порівняльна характеристика розглянутих рішень за ключовими
характеристиками наведена у табл.~\ref{tab:solutions}.

\begin{table}[H]
  \caption{Порівняння існуючих рішень для обчислення ШПФ}
  \label{tab:solutions}
  \centering
  \small
  \begin{tabular}{p{2.4cm}p{2.0cm}p{1.8cm}p{1.6cm}p{1.8cm}p{2.4cm}}
    \toprule
    Рішення & Тип & Платформа & Динамічна візуалізація & Відкритий доступ & Ключова особливість \\
    \midrule
    FFTW 3.x & SW бібліотека & CPU (x86, ARM) & Немає & Так & Адаптивний planner, cache-oblivious рекурсія~\cite{frigo2005} \\[2pt]
    ARM CMSIS-DSP & SW бібліотека & ARM Cortex-M & Немає & Так & Фіксований radix-4, SIMD-оптимізація~\cite{cmsis_dsp} \\[2pt]
    NVIDIA CuFFT & SW бібліотека & GPU (NVIDIA) & Немає & Ні & Масовий паралелізм, вендор-лок~\cite{du2012} \\[2pt]
    Nukada autotuner & SW + автотюнінг & GPU (CUDA) & Немає & Част. & Runtime-компіляція~\cite{nukada2009} \\[2pt]
    SPIRAL & Генератор коду & CPU / FPGA & Немає & Так & Пошук у просторі факторизацій SPL~\cite{puschel2004} \\[2pt]
    Intel IPP & SW бібліотека & Intel x86 & Немає & Ні & Примітиви AVX/SIMD~\cite{intel_ipp} \\[2pt]
    MATLAB Simulink (SDF) & Середовище моделювання & Desktop & Часткова & Ні & SDF-модель, комерційна ліцензія \\[2pt]
    draw.io / Visio & Діаграмний редактор & Desktop / Web & Ні (статика) & Част. & Лише статична схема DFG \\[2pt]
    \textbf{Симулятор (ця робота)} & Веб-застосунок & Браузер & \textbf{Так} & \textbf{Так} & \textbf{Інтерактивний DFG + числова верифікація} \\
    \bottomrule
  \end{tabular}
\end{table}

Як видно з таблиці, жодне з розглянутих рішень не поєднує одночасно
достовірного числового обчислення, інтерактивної візуалізації руху
токенів та веб-доступності без комерційної ліцензії. Саме ця
незаповнена ніша визначає предмет цієї роботи та обумовлює
формулювання задачі у підрозділі~1.5.


% ─────────────────────────────────────────────────
\subsection{Системний аналіз задачі}

Результати підрозд.~1.1--1.2 підтверджують наявність об'єктивного
розриву між математичною моделлю, формальним DFG/SDF-описом та
виробничими реалізаціями алгоритму ШПФ. Для структурованого опису
цього розриву та переходу до постановки задачі застосовується
системний аналіз: визначаються зацікавлені сторони, цілі розробки,
обмеження та критерії успіху.

\subsubsection*{Зацікавлені сторони (stakeholders)}

Потенційні користувачі та бенефіціари розроблюваного симулятора
згруповані у табл.~\ref{tab:stakeholders}.

\begin{table}[H]
  \caption{Зацікавлені сторони та їхні основні потреби}
  \label{tab:stakeholders}
  \centering
  \small
  \begin{tabular}{p{3.4cm}p{5.0cm}p{5.0cm}}
    \toprule
    Сторона & Потреба & Як симулятор її задовольняє \\
    \midrule
    Студенти спеціальності 122 «Комп'ютерні науки» & Зрозуміти внутрішню роботу алгоритму ШПФ, побачити паралелізм і токенну динаміку & Анімація руху токенів по DFG, крок-за-кроком режим, теплова карта активності вузлів \\[2pt]
    Викладачі курсів ЦОС та паралельних обчислень & Демонстраційний інструмент для лекційних і лабораторних занять & Веб-застосунок без встановлення, відтворюваність на будь-якій ОС \\[2pt]
    Дослідники dataflow-архітектур & Перевірка гіпотез щодо планування та критичного шляху графа ШПФ & Метрики throughput/latency/queue, event log, налаштовувані затримки \\[2pt]
    Інженери DSP/FPGA на ранній стадії проєктування & Попередня оцінка потокової моделі перед апаратною реалізацією & Параметричне налаштування \texttt{latency} і \texttt{delay}, експорт подій \\
    \bottomrule
  \end{tabular}
\end{table}

\subsubsection*{Основні цілі розробки}

На основі проаналізованої предметної області та потреб
зацікавлених сторін визначено ієрархію цілей.

\textbf{Генеральна ціль:} створити програмний симулятор, що
відтворює потокове виконання 16-точкового алгоритму ШПФ
radix-4 $4\times4$ та забезпечує його інтерактивну графічну
візуалізацію у вигляді анімованого DFG зі спостереженням токенів,
метрик і наскрізної затримки.

Підцілі, що випливають із генеральної:
\begin{enumerate}
  \item формалізувати граф потоку даних для алгоритму radix-4
        $4\times4$ з вузлами типів \texttt{source}, \texttt{dft4},
        \texttt{twiddle}, \texttt{sink} та дугами із параметричними
        затримками;
  \item реалізувати ядро симуляції з правилом активації токенів,
        FIFO-чергами та керованим логічним часом;
  \item забезпечити інтерактивну візуалізацію руху токенів,
        теплову карту активності вузлів і спектральний вихід;
  \item виконати числову верифікацію результатів шляхом порівняння
        з еталонним обчисленням ДПФ.
\end{enumerate}

\subsubsection*{Обмеження}

Розробка виконується у таких обмеженнях:

\begin{itemize}
  \item \textit{технічні} --- цільова платформа виконання: сучасний
        веб-браузер (Chrome/Firefox/Edge) з підтримкою JavaScript
        ES2020; відсутність серверної частини (повністю клієнтська
        симуляція);
  \item \textit{ресурсні} --- розробка силами одного виконавця
        у межах бакалаврської кваліфікаційної роботи;
  \item \textit{часові} --- термін виконання згідно з календарним
        графіком кафедри;
  \item \textit{алгоритмічні} --- фіксований розмір перетворення
        $N = 16$ як базовий демонстраційний розмір; розширення на
        довільні $N = 4^t$ залишається напрямком подальшого розвитку;
  \item \textit{правові} --- використання лише програмних бібліотек
        з відкритими ліцензіями (MIT, Apache~2.0, BSD) без
        комерційних компонентів.
\end{itemize}

\subsubsection*{Критерії успіху}

Результат роботи вважається досягнутим, якщо одночасно виконуються
такі вимірювані критерії:

\begin{itemize}
  \item максимальна відхилення вихідного спектра симулятора від
        еталонного ДПФ за модулем не перевищує $\varepsilon = 10^{-9}$;
  \item чотири вузли \texttt{dft4} першого ярусу виконуються
        паралельно, що підтверджується одночасними подіями
        \texttt{onFire} у журналі симуляції;
  \item користувач може керувати швидкістю симуляції у діапазоні
        не менш ніж $0.1\times$--$10\times$ без втрати коректності;
  \item час відгуку інтерфейсу на користувацькі дії (пуск, пауза,
        зміна швидкості) не перевищує 100~мс.
\end{itemize}


% ─────────────────────────────────────────────────
\subsection{Етичні, правові та регуляторні аспекти}

\subsubsection*{Характер даних, що опрацьовуються}

Розроблюваний симулятор працює з синтетичними тестовими сигналами
(синусоїдами, імпульсами, шумом), що генеруються безпосередньо у
браузері користувача. Проєкт \emph{не опрацьовує} персональних,
медичних, біометричних або фінансових даних, не передає жодних
даних на зовнішні сервери та не використовує моделей машинного
навчання для ухвалення рішень. Унаслідок цього більшість вимог
Загального регламенту захисту даних ЄС (GDPR) та Закону України
«Про захист персональних даних» (№~2297-VI) не поширюються на
продукт у прямій формі. Водночас базові принципи мінімізації даних
і прозорості обробки враховані: застосунок не веде телеметрії,
не встановлює ідентифікаційних cookies та працює повністю
на стороні клієнта.

\subsubsection*{Академічна доброчесність}

Реалізація симулятора здійснюється з дотриманням політики
академічної доброчесності Національного університету «Львівська
політехніка». Усі використані зовнішні джерела (публікації,
бібліотеки, референтні реалізації) мають явні посилання у списку
використаних джерел. Обчислювальне ядро симулятора написане
самостійно; еталонна бібліотека \texttt{fft.js} використовується
виключно як референтна для верифікації (підрозд.~4.2) і не
залучається до шляху виконання основного алгоритму. Плагіат,
фальсифікація та фабрикація результатів не допускаються.

\subsubsection*{Ліцензійні та правові вимоги до програмного забезпечення}

Робота базується на програмних компонентах із відкритими ліцензіями:
Next.js (MIT), React (MIT), React Flow (MIT), Zustand (MIT),
TypeScript (Apache~2.0). Такі ліцензії допускають некомерційне
навчальне використання, модифікацію та розповсюдження вихідного
коду за умови збереження оригінальних повідомлень про авторство.
Розроблений у цій роботі вихідний код оприлюднюється під відкритою
ліцензією, що забезпечує відтворюваність результатів та можливість
їх повторного використання у подальших дослідженнях і навчальних
курсах.

\subsubsection*{Відповідність галузевим стандартам}

При розробці враховуються такі чинні стандарти та технічні
регламенти:

\begin{itemize}
  \item \textbf{ISO/IEC 25010:2011} --- модель якості програмного
        забезпечення; для симулятора релевантними є характеристики
        \textit{функціональна придатність}, \textit{продуктивність}
        (час відгуку $<100$~мс), \textit{зручність використання}
        та \textit{надійність} (детермінованість SDF-графа);
  \item \textbf{IEEE 754} --- стандарт подання чисел з плаваючою
        крапкою; визначає точність обчислення комплексних операцій
        у подвійній точності (double, $\varepsilon_{\mathrm{mach}}
        \approx 2.22 \cdot 10^{-16}$);
  \item \textbf{ECMAScript 2020 (ES11)} --- стандартизація мови
        JavaScript, що забезпечує переносимість коду між сучасними
        браузерами;
  \item \textbf{ДСТУ 3008:2015} --- вимоги до оформлення звітів у
        сфері науки і техніки, що застосовуються до пояснювальної
        записки БКР;
  \item \textbf{ДСТУ ГОСТ 7.1:2006} --- правила бібліографічного
        опису, що застосовуються у списку використаних джерел.
\end{itemize}

Оскільки проєкт не використовує персональних даних і не залучає
моделей ШІ для ухвалення рішень, питання упередженості моделей
(bias/fairness), галюцинацій LLM та конфіденційності запитів у
межах цієї роботи є нерелевантними.


% ─────────────────────────────────────────────────
\subsection{Постановка задачі та формулювання вимог}

\subsubsection*{Формальна постановка задачі}

Нехай задано комплекснозначну вхідну послідовність
$\{x_n\}_{n=0}^{N-1}$ довжиною $N = 16$, подану у подвійній
точності формату IEEE~754. Потрібно розробити програмний
симулятор, який:

\begin{enumerate}
  \item обчислює дискретне перетворення Фур'є
        \begin{equation*}
          X_k = \sum_{n=0}^{N-1} x_n \, e^{-j 2\pi n k / N},
          \quad k = 0, 1, \ldots, N-1,
        \end{equation*}
        за алгоритмом radix-4 з двовимірною $4 \times 4$
        декомпозицією, реалізованим у моделі потоку даних;
  \item формалізує процес обчислення у вигляді орієнтованого
        графа потоку даних $G = (V, E)$, де $V$ --- множина вузлів
        типів \{\texttt{source}, \texttt{dft4}, \texttt{twiddle},
        \texttt{sink}\}, $E$ --- множина спрямованих дуг із
        параметричними затримками;
  \item виконує симуляцію графа з дотриманням правила активації
        (вузол спрацьовує тоді і лише тоді, коли на всіх його
        вхідних каналах є токени) та з урахуванням індивідуальних
        затримок \texttt{latency} вузлів і \texttt{delay} дуг;
  \item забезпечує інтерактивну візуалізацію процесу обчислення
        та подає кінцевий спектр $\{X_k\}$ у вигляді амплітудної
        діаграми.
\end{enumerate}

\subsubsection*{Вхідні та вихідні дані}

\textbf{Вхідні дані симулятора:}
\begin{itemize}
  \item послідовність $\{x_n\}_{n=0}^{15}$ --- 16~комплексних
        відліків сигналу, задається користувачем або генерується
        вбудованим джерелом (синусоїда, одиничний імпульс,
        білий шум);
  \item параметри симуляції --- затримки вузлів
        ($L_{\mathrm{dft4}}$, $L_{\mathrm{twiddle}}$), затримка
        дуг ($\Delta_{\mathrm{edge}}$), множник швидкості
        (speed multiplier);
  \item команди керування --- пуск, пауза, покроковий режим,
        скидання стану.
\end{itemize}

\textbf{Вихідні дані симулятора:}
\begin{itemize}
  \item спектр $\{X_k\}_{k=0}^{15}$ --- 16~комплексних відліків,
        що візуалізуються як амплітуда $|X_k|$ і фаза $\arg X_k$;
  \item журнал подій (event log) --- часові мітки подій
        \texttt{onToken}, \texttt{onFire}, \texttt{onOutput} для
        кожного вузла;
  \item метрики роботи графа --- пропускна здатність (throughput),
        наскрізна затримка (latency), розмір черг (queue size),
        теплова карта активності вузлів;
  \item звіт про верифікацію --- максимальне відхилення $\max_k
        |X_k^{\mathrm{sim}} - X_k^{\mathrm{ref}}|$ відносно
        еталонного ДПФ.
\end{itemize}

\subsubsection*{Функціональні вимоги}

Функціональні вимоги до симулятора зведено у табл.~\ref{tab:func_req}.

\begin{table}[H]
  \caption{Функціональні вимоги до симулятора}
  \label{tab:func_req}
  \centering
  \small
  \begin{tabular}{p{1.4cm}p{4.5cm}p{7.5cm}}
    \toprule
    ID & Назва & Опис \\
    \midrule
    F1 & Побудова DFG & Автоматична генерація графа потоку даних для ШПФ radix-4 $4\times4$ при $N = 16$ \\
    F2 & Симуляція виконання & Крокування часу з доставкою токенів, перевіркою правила активації та виконанням вузлів \\
    F3 & Введення сигналу & Вибір типу вхідного сигналу (синусоїда, імпульс, шум) та його параметрів \\
    F4 & Керування симуляцією & Пуск, пауза, покроковий режим, скидання, зміна множника швидкості \\
    F5 & Візуалізація токенів & Анімоване переміщення токенів по дугах з відображенням значень \\
    F6 & Відображення метрик & Показ throughput, latency, queue size, теплової карти активності \\
    F7 & Спектральний вихід & Відображення амплітуди та фази $\{X_k\}$ у вигляді гістограми \\
    F8 & Верифікація результату & Порівняння з еталонним ДПФ і виведення максимального відхилення \\
    F9 & Налаштування параметрів & Редагування затримок \texttt{latency}, \texttt{delay} без перезавантаження сторінки \\
    \bottomrule
  \end{tabular}
\end{table}

\subsubsection*{Нефункціональні вимоги}

Нефункціональні вимоги регламентують якісні характеристики системи
(табл.~\ref{tab:nonfunc_req}) у відповідності до ISO/IEC~25010.

\begin{table}[H]
  \caption{Нефункціональні вимоги до симулятора}
  \label{tab:nonfunc_req}
  \centering
  \small
  \begin{tabular}{p{1.4cm}p{3.8cm}p{8.2cm}}
    \toprule
    ID & Характеристика & Вимога \\
    \midrule
    NF1 & Точність обчислень & $\max_k |X_k^{\mathrm{sim}} - X_k^{\mathrm{ref}}| \le 10^{-9}$ \\
    NF2 & Час відгуку UI & Реакція на користувацьку дію $\le 100$~мс \\
    NF3 & Частота кадрів & Стабільні 60~fps при швидкості симуляції $\le 5\times$ \\
    NF4 & Портативність & Запуск у Chrome~$\ge$~90, Firefox~$\ge$~88, Edge~$\ge$~90 \\
    NF5 & Детермінованість & Однаковий вхід і параметри дають ідентичний журнал подій \\
    NF6 & Відкритий код & Розповсюдження під ліцензією MIT або Apache~2.0 \\
    NF7 & Без встановлення & Робота у браузері без серверної частини та БД \\
    NF8 & Зручність & Основні операції доступні не глибше двох кліків від головного екрану \\
    \bottomrule
  \end{tabular}
\end{table}

\subsubsection*{Діаграма прецедентів (Use Case)}

У системі виділяються два актори --- \emph{Студент} (первинний
користувач, спостерігач) і \emph{Викладач} (вторинний користувач,
може налаштовувати сценарії демонстрації). Основні прецеденти та
їхні зв'язки наведено у табл.~\ref{tab:usecases}.

\begin{table}[H]
  \caption{Перелік прецедентів (Use Cases) симулятора}
  \label{tab:usecases}
  \centering
  \small
  \begin{tabular}{p{1.2cm}p{4.0cm}p{2.0cm}p{6.5cm}}
    \toprule
    ID & Прецедент & Актор & Стислий опис \\
    \midrule
    UC1 & Завантажити готовий сценарій & Студент, Викладач & Обрати один із передустановлених входів (синусоїда, імпульс, шум) \\
    UC2 & Задати власний вхідний сигнал & Студент & Ввести 16~комплексних відліків вручну \\
    UC3 & Запустити симуляцію & Студент & Ініціювати виконання DFG з пуску до завершення \\
    UC4 & Керувати темпом симуляції & Студент & Пауза, step-by-step, зміна множника швидкості \\
    UC5 & Спостерігати рух токенів & Студент & Переглядати анімацію токенів по дугах у реальному часі \\
    UC6 & Переглянути спектр результату & Студент & Побачити $|X_k|$ і $\arg X_k$ у вигляді діаграми \\
    UC7 & Порівняти з еталоном & Студент, Викладач & Запустити \texttt{compareWithFFT()} і побачити похибку \\
    UC8 & Змінити параметри затримок & Викладач & Налаштувати \texttt{latency} вузлів і \texttt{delay} дуг \\
    UC9 & Експортувати журнал подій & Викладач & Зберегти event log у файл для аналізу \\
    \bottomrule
  \end{tabular}
\end{table}

\emph{Підпорядкування прецедентів:} UC3 містить (\textit{include})
UC5 (рух токенів автоматично супроводжує виконання); UC7
розширює (\textit{extend}) UC3 та може бути викликаний після
завершення симуляції; UC1 і UC2 є взаємовиключними альтернативами
підготовки вхідних даних.

Описана постановка задачі та перелік вимог утворюють технічне
завдання, яке далі конкретизується у розділі~2 через вибір методів
і технологій реалізації, проєктування архітектури та структур
даних.


% ─────────────────────────────────────────────────
\subsection*{Висновки до розділу~1}

У першому розділі виконано аналіз предметної області, огляд сучасних
підходів, системний аналіз задачі, розгляд етичних та правових
аспектів і формалізовано постановку задачі.

Зафіксовано \emph{стан предметної області:} алгоритм ШПФ radix-4
$4\times4$ для $N = 16$ має достатньо розроблений математичний
апарат і формальні потокові моделі (DFG, SDF), однак у практичному
інструментарії ці рівні описані відокремлено. Виробничі бібліотеки
(FFTW, ARM~CMSIS-DSP, NVIDIA~CuFFT, SPIRAL, Intel~IPP) оптимізовані
виключно для продуктивності та не надають засобів інтроспекції.
Методи візуального моделювання (UML, IDEF0, DFG/SDF, сигнальні
графи) залишаються переважно теоретичними та не реалізуються як
інтерактивні веб-симулятори з числовою верифікацією.

Виявлено \emph{прогалину:} жодне з розглянутих рішень не поєднує
одночасно достовірного обчислення ШПФ за формальною SDF-моделлю,
інтерактивної динамічної візуалізації руху токенів і веб-доступності
без комерційної ліцензії. Саме ця незаповнена ніша визначає предмет
цієї роботи.

Сформульовано \emph{задачу} розробки програмного симулятора потокового
обчислювача 16-точкового ШПФ з анімованим DFG, визначено вхідні та
вихідні дані, технічні, ресурсні й правові обмеження та критерії
успіху. Розроблено перелік із 9~функціональних (F1--F9) і 8~нефункціональних
(NF1--NF8) вимог, побудовано перелік із 9~прецедентів (UC1--UC9),
що охоплюють основні сценарії використання симулятора студентами
та викладачами. Ключові вимірювані критерії --- точність $\varepsilon
\le 10^{-9}$, час відгуку $\le 100$~мс, стабільні 60~fps при
швидкості до $5\times$ --- підлягають експериментальній перевірці
у розділі~4.

Розглянуто етичні та правові аспекти: робота не опрацьовує
персональних даних, базується на відкритих ліцензіях (MIT,
Apache~2.0) і відповідає стандартам ISO/IEC~25010, IEEE~754,
ECMAScript~2020, ДСТУ~3008:2015 та ДСТУ~ГОСТ~7.1:2006.

Математичне та алгоритмічне обґрунтування (виведення формул
radix-4 $4\times4$, матричні факторизації, формалізація DFG/SDF)
виноситься до розділу~2, де воно подається у складі проєктного
обґрунтування обраних методів і архітектури системи.

\newpage

% --- Розділ 2 ---
% chapters/chapter2.tex
\section{ПРОЄКТУВАННЯ ПРОГРАМНОГО РІШЕННЯ}

% Розділ 2 згідно нових методичних вказівок (12-15 стор.):
% 2.1 Вибір методів, технологій та засобів реалізації
% 2.2 Математичне та алгоритмічне забезпечення
% 2.3 Загальна архітектура / структура системи
% 2.4 Структури даних
% 2.5 Механізми обміну даними та інтеграції системи
% 2.6 Інтерфейс користувача та взаємодія з системою
% 2.7 Висновки до розділу 2
%
% Конкретні версії бібліотек, середовища розробки та деталі реалізації
% виносяться до розділу 3.

% ─────────────────────────────────────────────────
\subsection{Вибір методів, технологій та засобів реалізації}

У підрозд.~1.5 сформульовано технічне завдання: розробити веб-симулятор
потокового обчислювача 16-точкового ШПФ з інтерактивною візуалізацією
руху токенів, числовою верифікацією та відкритим доступом. Нижче
обґрунтовується вибір архітектурних підходів і сімейств інструментів,
що покривають ці вимоги. Конкретні версії бібліотек та параметри
середовища розробки наводяться у підрозд.~3.1.

\subsubsection*{Вибір архітектурного підходу}

Розглянуто три альтернативних підходи до побудови симулятора.

\textbf{Варіант~А: настільний застосунок (Electron, Qt, JavaFX).}
Забезпечує повний доступ до ресурсів хост-системи та багатопоточність,
однак вимагає встановлення на кожну робочу станцію, різних збірок під
Windows, Linux, macOS та погіршує переносимість --- прямо суперечить
вимозі NF7 (робота без встановлення) і ускладнює дотримання NF4
(портативність).

\textbf{Варіант~Б: клієнт-серверний веб-застосунок.} Тонкий клієнт у
браузері виконує лише візуалізацію, а ядро симуляції розміщується на
сервері. Такий підхід вносить мережеву затримку, що несумісно з
вимогою NF2 (час відгуку $\le 100$~мс), потребує підтримки серверної
інфраструктури та ускладнює відтворюваність результатів поза межами
хостингу.

\textbf{Варіант~В: повністю клієнтський SPA (Single-Page Application).}
--- \emph{обрано як основу архітектури}. Усі обчислення виконуються
на стороні браузера; серверна частина відсутня як така. Завантаження
застосунку відбувається одноразово як статичний пакет (HTML+JS+CSS),
після чого робота не потребує мережі. Цей варіант одночасно задовольняє
NF4, NF5 (детермінованість локального виконання), NF7 (без встановлення)
і забезпечує низький час відгуку за рахунок відсутності мережевих
викликів у петлі симуляції.

\subsubsection*{Вибір парадигми реалізації UI та стану}

Симулятор поєднує три принципово різні підсистеми: анімована
графова візуалізація DFG, числове ядро симуляції та панель керування
і метрик. Для їх ефективної координації обрано поєднання
\emph{реактивного рендерингу} та \emph{централізованого
спостерігаваного сховища стану}.

Реактивний рендеринг (React-парадигма) природно відображає
стан симулятора на візуальні компоненти: зміна значення в одному
полі сховища автоматично ініціює перемалювання лише тих компонентів,
що підписані на це поле. Це усуває ручне управління DOM і вкладається
у вимогу NF3 (стабільні 60~fps при швидкості симуляції $\le 5\times$).

Централізоване сховище стану (pattern single source of truth)
гарантує, що всі панелі --- GraphView, SpectrumView, Metrics ---
бачать один і той самий «знімок» стану у поточному кадрі. Це
виключає класичні помилки неузгодженості між візуальними
представленнями однієї сутності.

Альтернатива у вигляді прямого маніпулювання DOM (jQuery-style)
була відкинута через лінійне зростання складності при додаванні
нових компонентів, а Flux-підходи з ред'юсерами (Redux) --- через
надмірну кількість шаблонного коду для базових сценаріїв.

\subsubsection*{Вибір методу візуалізації DFG}

Для реалізації інтерактивного графа потоку даних розглянуто три
підходи:

\begin{itemize}
  \item \textbf{Рендеринг на HTML-canvas.} Забезпечує максимальну
        продуктивність при великій кількості об'єктів, але вимагає
        ручної реалізації перетягування, виділення та спливаючих
        підказок, що значно збільшує обсяг коду для UI-інтеракцій.
  \item \textbf{SVG + ручне кодування.} Простіший у налагодженні,
        але лише для невеликих графів; розширюваність обмежена
        тим, що кожен новий тип вузла потребує окремого обробника
        подій.
  \item \textbf{Спеціалізована бібліотека графової візуалізації
        (React~Flow).} --- \emph{обрано}. Надає готові механізми
        перетягування вузлів, зума, анімованих ребер, плаваючих
        підказок, а також формальну модель \texttt{Node}/\texttt{Edge},
        що добре відображається на абстракцію DFG.
\end{itemize}

\subsubsection*{Шаблони проєктування}

У проєктуванні системи застосовано такі класичні шаблони:

\begin{itemize}
  \item \textit{Observer (Publisher/Subscriber)} --- для подій
        \texttt{onToken}, \texttt{onFire}, \texttt{onOutput}:
        ядро симуляції є видавцем, а UI-компоненти
        (Metrics, теплова карта, SpectrumView) --- підписниками
        без прямих посилань на ядро;
  \item \textit{State / Finite State Machine} --- для керування
        життєвим циклом симуляції: \texttt{IDLE} $\to$
        \texttt{RUNNING} $\to$ \texttt{PAUSED} $\to$ \texttt{DONE};
  \item \textit{Strategy} --- для реалізації різних типів вузлів
        (\texttt{source}, \texttt{dft4}, \texttt{twiddle},
        \texttt{sink}): поведінка вузла інкапсулюється у функції,
        обрану за полем \texttt{NodeKind};
  \item \textit{Factory} --- для генерації графа через функцію
        \texttt{generateDFT4x4()}, що повертає готовий об'єкт
        \texttt{Graph} на основі параметрів $N$ та \texttt{latency};
  \item \textit{Facade} --- клас \texttt{DataflowRuntime} приховує
        внутрішню реалізацію черг і планувальника за єдиним методом
        \texttt{tick(dt, maxFires)}.
\end{itemize}

\subsubsection*{Типізація та засоби забезпечення якості}

Для дотримання NF1 (точність $\varepsilon \le 10^{-9}$) і NF5
(детермінованість) проєктом передбачено статичну типізацію вихідного
коду. Це унеможливлює помилки передачі аргументів неправильного
типу (наприклад, числа замість комплексного об'єкта) на етапі
компіляції. Комплексні числа подаються як структура \{\,re, im\,\}
з операціями у подвійній точності IEEE~754, що забезпечує необхідний
запас точності відносно $\varepsilon$.

Підсумкова таблиця зв'язків між вимогами підрозд.~1.5 та обраними
архітектурними рішеннями наведена у табл.~\ref{tab:req_to_tech}.

\begin{table}[H]
  \caption{Зв'язок вимог і архітектурних рішень}
  \label{tab:req_to_tech}
  \centering
  \small
  \begin{tabular}{p{2.2cm}p{6.0cm}p{5.8cm}}
    \toprule
    Вимога & Архітектурне рішення & Обґрунтування \\
    \midrule
    NF1 (точність) & Статична типізація + IEEE 754 double & Виключення помилок типів, $\varepsilon_{\mathrm{mach}}{\approx}2.22{\cdot}10^{-16}$ \\
    NF2 (відгук $\le 100$~мс) & Клієнтський SPA без серверних викликів & Відсутність мережевої затримки \\
    NF3 (60~fps) & Реактивний рендеринг з підпискою на окремі поля & Перемальовування лише змінених компонентів \\
    NF4 (портативність) & Стандарт ES2020, Web APIs & Робота у Chrome/Firefox/Edge без polyfill \\
    NF5 (детермінованість) & Логічний час + статичний граф & Ідентичні входи $\Rightarrow$ ідентичний журнал \\
    NF6 (MIT/Apache) & Вибір лише OSS-бібліотек & React, React Flow, Zustand, TypeScript \\
    NF7 (без встановлення) & SPA у статичному хостингу & Одноразове завантаження бандлу \\
    NF8 (зручність) & React Flow + панелі Controls/Metrics & Двокліковий доступ до базових операцій \\
    \bottomrule
  \end{tabular}
\end{table}


% ─────────────────────────────────────────────────
\subsection{Математичне та алгоритмічне забезпечення}

У цьому підрозділі подано математичний апарат 16-точкового ШПФ
з основою~4, що становить інтелектуальне ядро симулятора, та формалізацію
обчислень у парадигмі потоку даних. Виведені формули та схеми алгоритму
безпосередньо визначають структуру графа (підрозд.~2.3) і структури
даних (підрозд.~2.4).

\subsubsection*{Дискретне перетворення Фур'є}

Дискретне перетворення Фур'є (ДПФ) скінченної комплекснозначної
послідовності $\{x_n\}_{n=0}^{N-1}$ визначається
як~\cite{nussbaumer1982,protsko2025}:

\begin{equation}
  X_k = \sum_{n=0}^{N-1} x_n \, e^{-j 2\pi n k / N}
       = \sum_{n=0}^{N-1} x_n \, W_N^{n k},
  \quad k = 0, 1, \ldots, N-1,
  \label{eq:dft}
\end{equation}

де $W_N = e^{-j 2\pi/N}$ --- поворотний множник (twiddle factor).
Пряме обчислення~(\ref{eq:dft}) потребує $N^2$ комплексних множень
і $N(N-1)$ комплексних додавань, тобто має складність $O(N^2)$, що
недостатньо для задач реального часу при великих $N$.

Клас швидких алгоритмів перетворення Фур'є (ШПФ) знижує складність
до $O(N \log_2 N)$ за рахунок факторизації довжини
$N = N_1 \cdot N_2$ та заміни одного $N$-точкового ДПФ на систему
$N_1$ коротших $N_2$-точкових та $N_2$ коротших $N_1$-точкових
перетворень~\cite{cooley1965,nussbaumer1982}. Для $N = 16$ природним
вибором є алгоритм з основою~4 (radix-4), у якому $N_1 = N_2 = 4$.

\subsubsection*{Двовимірна $4 \times 4$ декомпозиція для $N = 16$}

Для $N = 16$ часовий і частотний індекси подаються у вигляді:

\begin{equation}
  n = 4s + r, \quad k = 4k_2 + k_1, \quad r, s, k_1, k_2 \in \{0,1,2,3\}.
  \label{eq:idx_map}
\end{equation}

Вхідна послідовність $\{x(n)\}_{n=0}^{15}$ вкладається у матрицю
$\mathbf{A}$ розміром $4 \times 4$:

\begin{equation}
  A[r][s] = x(4s + r), \quad r, s \in \{0,1,2,3\}.
  \label{eq:mat_A}
\end{equation}

Підстановка~(\ref{eq:idx_map}) у~(\ref{eq:dft}) із використанням
тотожності $W_{16}^{4 \cdot (\cdot)} = W_4^{(\cdot)}$
дає~\cite{burrus1985dft,nussbaumer1982}:

\begin{equation}
  X(4k_2 + k_1)
    = \sum_{r=0}^{3} W_{16}^{r k_1}
      \underbrace{\Bigl( \sum_{s=0}^{3} A[r][s]\, W_4^{s k_1} \Bigr)}_{G_r(k_1)}
      W_4^{r k_2},
  \label{eq:fft4x4}
\end{equation}

що визначає тристадійну схему алгоритму.

\textbf{Стадія~1 --- рядкові 4-точкові ДПФ.}
Для кожного $r \in \{0,1,2,3\}$ обчислюється 4-точкове ДПФ рядка:

\begin{equation}
  G_r(k_1) = \sum_{s=0}^{3} A[r][s]\, W_4^{s k_1},
  \quad k_1 \in \{0,1,2,3\}.
  \label{eq:stage1}
\end{equation}

Чотири перетворення~(\ref{eq:stage1}) є взаємно незалежними та
виконуються паралельно~\cite{nussbaumer1982}.

\textbf{Стадія~2 --- поворотні множники.}
Результати стадії~1 множаться на матрицю поворотних множників:

\begin{equation}
  H_r(k_1) = W_{16}^{r \cdot k_1} \cdot G_r(k_1),
  \quad r, k_1 \in \{0,1,2,3\}.
  \label{eq:stage2}
\end{equation}

При $r = 0$ або $k_1 = 0$ маємо $W_{16}^0 = 1$, тому 7~з 16~множень
є тривіальними. Числові значення 9~нетривіальних множників наведено
у табл.~\ref{tab:twiddle}.

\begin{table}[H]
  \caption{Поворотні множники $W_{16}^{r \cdot k_1}$ (округлено до 3 знаків)}
  \label{tab:twiddle}
  \centering
  \small
  \begin{tabular}{c|cccc}
    \toprule
    $r \backslash k_1$ & 0 & 1 & 2 & 3 \\
    \midrule
    0 & $1$ & $1$ & $1$ & $1$ \\
    1 & $1$ & $0{,}924{-}0{,}383j$ & $0{,}707{-}0{,}707j$ & $0{,}383{-}0{,}924j$ \\
    2 & $1$ & $0{,}707{-}0{,}707j$ & $-j$                  & $-0{,}707{-}0{,}707j$ \\
    3 & $1$ & $0{,}383{-}0{,}924j$ & $-0{,}707{-}0{,}707j$ & $-0{,}924{+}0{,}383j$ \\
    \bottomrule
  \end{tabular}
\end{table}

\textbf{Стадія~3 --- стовпцеві 4-точкові ДПФ.}
Для кожного $k_1$ обчислюється 4-точкове ДПФ стовпця матриці
$\mathbf{H}$:

\begin{equation}
  X(4k_1 + k_2) = \sum_{r=0}^{3} H_r(k_1)\, W_4^{r k_2},
  \quad k_2 \in \{0,1,2,3\}.
  \label{eq:stage3}
\end{equation}

Чотири стовпцевих перетворення~(\ref{eq:stage3}) також є взаємно
незалежними.

\subsubsection*{Оптимізоване 4-точкове ДПФ (метелик radix-4)}

Ядром обох стадій~(\ref{eq:stage1}) і~(\ref{eq:stage3}) є 4-точкове
ДПФ, що завдяки властивостям $W_4 \in \{1, -j, -1, j\}$ реалізується
без нетривіальних множень. Для входів $v_0, v_1, v_2, v_3 \in \mathbb{C}$
вводяться проміжні суми/різниці:

\begin{equation*}
  E_0 = v_0 + v_2, \quad E_1 = v_0 - v_2,
\end{equation*}
\begin{equation*}
  O_0 = v_1 + v_3, \quad O_1 = v_1 - v_3,
\end{equation*}

після чого вихідні відліки отримуються як:

\begin{equation}
  \begin{aligned}
    V_0 &= E_0 + O_0, & V_2 &= E_0 - O_0, \\
    V_1 &= E_1 - j \cdot O_1, & V_3 &= E_1 + j \cdot O_1.
  \end{aligned}
  \label{eq:dft4_butterfly}
\end{equation}

Множення на $j$ у~(\ref{eq:dft4_butterfly}) зводиться до перестановки
дійсної та уявної частин зі зміною знаку
($j \cdot (a + bj) = -b + aj$) і не потребує операції множення у
традиційному сенсі. Таким чином, одне 4-точкове ДПФ потребує
\textbf{8~комплексних додавань і 0 нетривіальних
множень}~\cite{nussbaumer1982}.

\subsubsection*{Підрахунок операцій}

Загальна обчислювальна вартість алгоритму для $N = 16$:

\begin{itemize}
  \item \textbf{Стадія~1:} $4 \cdot 8 = 32$ комплексних додавань, $0$
        нетривіальних множень;
  \item \textbf{Стадія~2:} $9$ нетривіальних комплексних множень
        (що відповідає $36$ дійсним множенням), $16$ комплексних
        додавань;
  \item \textbf{Стадія~3:} $4 \cdot 8 = 32$ комплексних додавань, $0$
        нетривіальних множень.
\end{itemize}

Разом: \textbf{9 нетривіальних комплексних множень} і
\textbf{80 комплексних додавань}, порівняно з $N^2 = 256$ комплексних
множень прямого ДПФ~(\ref{eq:dft}) --- прискорення у
$\approx 28$~разів за кількістю множень~\cite{nussbaumer1982,burrus1985dft}.

\subsubsection*{Формалізація у парадигмі потоку даних}

Алгоритм~(\ref{eq:stage1})--(\ref{eq:stage3}) подається у вигляді
орієнтованого ациклічного графа потоку даних (DFG)
$G = (V, E)$~\cite{dennis1974,veen1986,lee1994}, де:

\begin{itemize}
  \item $V$ --- множина вузлів чотирьох типів
        (\texttt{source}, \texttt{dft4}, \texttt{twiddle}, \texttt{sink});
  \item $E \subseteq V \times V$ --- множина спрямованих дуг, що
        є FIFO-чергами токенів.
\end{itemize}

\textbf{Правило активації.} Вузол $v \in V$ є готовим до виконання
тоді і лише тоді, коли на кожному з його вхідних каналів присутній
щонайменше один токен. При активації вузол поглинає вхідні токени,
виконує обчислення, визначене його типом, та породжує відповідні
вихідні токени, що розміщуються на вихідних дугах.

\textbf{Часова модель.} Вузол має затримку виконання \texttt{latency},
а дуга --- затримку доставки \texttt{delay}. Токен, породжений у
момент $t$, стає готовим до споживання у момент $t + \texttt{delay}$.
Логічний час $\tau$ симулятора інкрементується дискретно на кроці
$\mathrm{d}\tau$, що регулюється множником швидкості.

\subsubsection*{Псевдокод планувальника}

Планувальник (ядро \texttt{DataflowRuntime}) на кожному кроці
симуляції виконує дві дії: доставку токенів із FIFO-черг і запуск
готових вузлів. Узагальнений псевдокод:

\begin{figure}[H]
\begin{verbatim}
procedure TICK(dt, maxFires)
    tau <- tau + dt
    for each edge e in E
        while e.queue.head.t + e.delay <= tau
            token <- e.queue.pop()
            v <- target(e); port <- targetPort(e)
            v.inputs[port] <- token
    fires <- 0
    for each node v in V  order by id
        if v is ready and not busy and fires < maxFires
            emit onFire(v, tau)
            v.busy <- tau + v.latency
            outs <- compute(v.kind, v.inputs)
            schedule outs on outgoing edges at time v.busy
            fires <- fires + 1
    for each node v in V
        if v.busy <= tau then v.busy <- 0
\end{verbatim}
\caption{Псевдокод однієї ітерації ядра симуляції \texttt{DataflowRuntime.tick()}.
         Джерело: розроблено автором.}
\label{fig:tick_pseudo}
\end{figure}

Ця процедура забезпечує детермінованість (NF5): при фіксованому порядку
обходу вузлів результат симуляції не залежить від реального часу
виконання у браузері.


% ─────────────────────────────────────────────────
\subsection{Загальна архітектура системи}

Архітектура симулятора подається на двох рівнях моделі C4~\cite{c4model}
--- Context і Container, --- доповнених діаграмою потоків даних (DFD)
для внутрішньої логіки симуляції.

\subsubsection*{Контекстний рівень (C4 Level 1)}

На контекстному рівні система представляється як «чорна скринька»,
що взаємодіє з одним актором --- \emph{користувачем} (студент або
викладач), та з двома зовнішніми сутностями: \emph{браузером} (runtime-
середовище) і \emph{локальним сховищем браузера} (для збереження
пресетів вхідних сигналів і параметрів). Застосунок завантажується
одноразово зі статичного веб-хостингу та надалі працює автономно.

\subsubsection*{Рівень контейнерів (C4 Level 2)}

На рівні контейнерів симулятор декомпозується на чотири логічні
підсистеми, що виконуються у межах одного процесу браузера:

\begin{enumerate}
  \item \textbf{Simulation Core} --- ядро симуляції (\texttt{DataflowRuntime},
        \texttt{scheduler}, \texttt{math}): реалізує псевдокод
        рис.~\ref{fig:tick_pseudo}, зберігає стан графа та черг;
  \item \textbf{Graph Generator} --- модуль побудови DFG для заданого
        $N$ та параметрів затримок;
  \item \textbf{State Store} --- централізоване сховище стану, до
        якого підписуються всі UI-компоненти;
  \item \textbf{UI Layer} --- набір React-компонентів:
        \texttt{GraphView} (полотно React Flow), \texttt{SpectrumView}
        (діаграма $|X(k)|$), \texttt{Metrics} (пропускна здатність,
        затримка, черги), \texttt{Controls} (пуск/пауза/швидкість).
\end{enumerate}

Взаємодії між контейнерами: UI Layer читає стан зі State Store та
викликає команди (пуск, пауза); State Store делегує обробку команд
до Simulation Core; Simulation Core публікує події
(\texttt{onToken}, \texttt{onFire}, \texttt{onOutput}), на які
State Store оновлює відповідні поля; Graph Generator викликається
один раз при ініціалізації та за кожної зміни $N$ чи latency.

\subsubsection*{Граф потоку даних обчислювача}

Детальна структура DFG для $N = 16$ складається з \textbf{56~вузлів}
і \textbf{64~дуг}, розподілених за п'ятьма ярусами:

\begin{enumerate}
  \item 16~вузлів \texttt{source} (src$_0 \ldots$src$_{15}$) ---
        генерують токени вхідних відліків $x(n)$;
  \item 4~вузли \texttt{dft4} (row-0 \ldots row-3) --- виконують
        стадію~1~(\ref{eq:stage1});
  \item 16~вузлів \texttt{twiddle} (tw-$r$-$k_1$) --- виконують
        стадію~2~(\ref{eq:stage2});
  \item 4~вузли \texttt{dft4} (col-0 \ldots col-3) --- виконують
        стадію~3~(\ref{eq:stage3});
  \item 16~вузлів \texttt{sink} (snk$_0 \ldots$snk$_{15}$) ---
        збирають спектральні відліки $X(k)$.
\end{enumerate}

Характеристики вузлів і дуг зведено у
табл.~\ref{tab:nodes}--\ref{tab:edges}.

\begin{table}[H]
  \caption{Типи вузлів DFG та їхні характеристики}
  \label{tab:nodes}
  \centering
  \small
  \begin{tabular}{lccccl}
    \toprule
    Тип вузла & К-сть & Вх. портів & Вих. портів & Затримка & Функція \\
    \midrule
    \texttt{source}  & 16 & 0 & 1 & ---  & Генерація $x(n)$ \\
    \texttt{dft4}    &  8 & 4 & 4 & 400  & 4-точкове ДПФ~(\ref{eq:dft4_butterfly}) \\
    \texttt{twiddle} & 16 & 1 & 1 & 200  & Множення на $W_{16}^{r k_1}$ \\
    \texttt{sink}    & 16 & 1 & 0 & ---  & Збір $X(k)$ \\
    \midrule
    \multicolumn{2}{l}{\textbf{Разом: 56 вузлів}} & & & & \\
    \bottomrule
  \end{tabular}
\end{table}

\begin{table}[H]
  \caption{Типи дуг DFG}
  \label{tab:edges}
  \centering
  \small
  \begin{tabular}{llcc}
    \toprule
    З вузла & До вузла & К-сть дуг & Затримка \\
    \midrule
    \texttt{source}      & \texttt{dft4} (рядковий)  & 16 & 600 \\
    \texttt{dft4} (ряд.) & \texttt{twiddle}          & 16 & 600 \\
    \texttt{twiddle}     & \texttt{dft4} (стовпц.)   & 16 & 600 \\
    \texttt{dft4} (ст.)  & \texttt{sink}             & 16 & 600 \\
    \midrule
    \multicolumn{2}{l}{\textbf{Разом: 64 дуги}} & & \\
    \bottomrule
  \end{tabular}
\end{table}

\subsubsection*{Критичний шлях і коефіцієнт паралелізму}

Критичний шлях у DFG (найдовший шлях від \texttt{source} до
\texttt{sink} з урахуванням затримок вузлів і дуг) становить:

\begin{equation}
  T_{\mathrm{kr}} = 400 + 600 + 200 + 600 + 400 = 2200
  \label{eq:critical_path}
\end{equation}

умовних одиниць часу. Послідовний час виконання (при одному
активному вузлі) становив би:

\begin{equation}
  T_{\mathrm{ps}} = 8 \cdot 400 + 16 \cdot 200 = 6400.
  \label{eq:seq_time}
\end{equation}

Коефіцієнт паралелізму:

\begin{equation}
  K_p = \frac{T_{\mathrm{ps}}}{T_{\mathrm{kr}}} = \frac{6400}{2200} \approx 2{,}9,
  \label{eq:parallelism}
\end{equation}

тобто паралельне виконання незалежних вузлів забезпечує майже
трикратне прискорення порівняно з послідовним обчисленням~\cite{veen1986}.

\begin{figure}[H]
  \centering
  \IfFileExists{images/dfg.png}{%
    \includegraphics[width=0.95\textwidth]{images/dfg.png}%
  }{%
    \fbox{\parbox[c][0.28\textheight][c]{0.95\textwidth}{\centering
      Рисунок буде додано після генерації у симуляторі.\\
      Файл: \texttt{tex/images/dfg.png}}}%
  }
  \caption{Граф потоку даних обчислювача 16-точкового ШПФ radix-4
           $4 \times 4$. Джерело: розроблено автором.}
  \label{fig:dfg}
\end{figure}


% ─────────────────────────────────────────────────
\subsection{Структури даних}

Інформаційне забезпечення симулятора визначається набором типів
даних, що є програмними відповідниками формальних об'єктів DFG
з підрозд.~2.2.

\subsubsection*{Комплексні числа}

Усі числові значення, з якими оперує симулятор, є комплексними.
Вони подаються структурою \texttt{Complex} у форматі
\{\,\texttt{re}: number,\ \texttt{im}: number\,\}, де \texttt{re},
\texttt{im} --- дійсна та уявна частини числа
$z = \texttt{re} + j \cdot \texttt{im}$, збережені у подвійній
точності IEEE~754 (64~біти, машинна точність
$\varepsilon_{\mathrm{mach}} \approx 2.22 \cdot 10^{-16}$).

\subsubsection*{Тип \texttt{Token}}

Тип \texttt{Token} є програмним відповідником токена DFG, що
переміщується по дугам графа (табл.~\ref{tab:token}).

\begin{table}[H]
  \caption{Поля типу \texttt{Token}}
  \label{tab:token}
  \centering
  \small
  \begin{tabular}{lll}
    \toprule
    Поле & Тип & Призначення \\
    \midrule
    \texttt{id}      & \texttt{string}               & Унікальний ідентифікатор токена \\
    \texttt{value}   & \texttt{Complex}              & Комплексне значення токена \\
    \texttt{t}       & \texttt{number}               & Логічний час надходження на дугу \\
    \texttt{originT} & \texttt{number}               & Час народження у \texttt{source} \\
    \bottomrule
  \end{tabular}
\end{table}

Поле \texttt{t} визначає момент готовності токена: він стає доступним
вузлу-приймачу в момент $t + \texttt{edge.delay}$. Поле
\texttt{originT} зберігається незмінним упродовж усього маршруту
токена і використовується для вимірювання наскрізної
затримки~(\ref{eq:critical_path}).

\subsubsection*{Тип \texttt{Edge}}

Тип \texttt{Edge} описує спрямований канал між двома портами
вузлів (табл.~\ref{tab:edge}).

\begin{table}[H]
  \caption{Поля типу \texttt{Edge}}
  \label{tab:edge}
  \centering
  \small
  \begin{tabular}{lll}
    \toprule
    Поле & Тип & Призначення \\
    \midrule
    \texttt{id}     & \texttt{string}                  & Унікальний ідентифікатор дуги \\
    \texttt{from}   & \texttt{\{node, port\}}          & Вузол і порт-джерело \\
    \texttt{to}     & \texttt{\{node, port\}}          & Вузол і порт-приймач \\
    \texttt{delay}  & \texttt{number}                  & Затримка доставки токена (600) \\
    \texttt{label}  & \texttt{string?}                 & Мітка на ребрі ($W_{16}^{rk_1}$) \\
    \texttt{labelValue} & \texttt{Complex?}            & Числове значення поворотного множника \\
    \bottomrule
  \end{tabular}
\end{table}

\subsubsection*{Тип \texttt{NodeSpec} і \texttt{Graph}}

Тип \texttt{NodeSpec} визначає специфікацію вузла DFG
(табл.~\ref{tab:nodespec}), а тип \texttt{Graph = \{nodes:
NodeSpec[], edges: Edge[]\}} --- повний граф обчислювача. Для
$N = 16$ об'єкт \texttt{Graph} містить 56~специфікацій вузлів та
64~дуги.

\begin{table}[H]
  \caption{Поля типу \texttt{NodeSpec}}
  \label{tab:nodespec}
  \centering
  \small
  \begin{tabular}{lll}
    \toprule
    Поле & Тип & Призначення \\
    \midrule
    \texttt{id}       & \texttt{string}    & Ідентифікатор вузла \\
    \texttt{kind}     & \texttt{NodeKind}  & Тип (source, dft4, twiddle, sink) \\
    \texttt{inPorts}  & \texttt{string[]}  & Список вхідних портів \\
    \texttt{outPorts} & \texttt{string[]}  & Список вихідних портів \\
    \texttt{latency}  & \texttt{number}    & Тривалість обчислення (400 або 200) \\
    \texttt{params}   & \texttt{object?}   & Параметри (для twiddle: $N, k$) \\
    \bottomrule
  \end{tabular}
\end{table}

Тип \texttt{NodeKind} є переліковим: \texttt{'source' | 'dft4' |
'twiddle' | 'sink'}. Він використовується як ключ у шаблоні
Strategy (підрозд.~2.1) для вибору функції обчислення вузла.

\subsubsection*{Проміжні матриці обчислення}

Алгоритм~(\ref{eq:stage1})--(\ref{eq:stage3}) оперує чотирма
матрицями даних розміром $4 \times 4$:

\begin{align}
  \mathbf{A}[r][s]   &= x(4s + r),            && \text{вхідні дані,}  \label{eq:A_def} \\
  \mathbf{G}[r][k_1] &= \mathrm{DFT4}(\mathbf{A}[r][\cdot])[k_1],
                     && \text{після стадії 1,}  \label{eq:G_def} \\
  \mathbf{H}[r][k_1] &= W_{16}^{r \cdot k_1} \cdot \mathbf{G}[r][k_1],
                     && \text{після стадії 2,}  \label{eq:H_def} \\
  X(4k_1 + k_2)      &= \mathrm{DFT4}(\mathbf{H}[\cdot][k_1])[k_2].
                     && \text{вихідний спектр.} \label{eq:X_def}
\end{align}

\subsubsection*{Числовий приклад: $x(n) = n$}

Для контрольної верифікації використано лінійний вхідний сигнал
$x(n) = n$, $n = 0 \ldots 15$. Проміжні значення для першого рядка
($r = 0$) і перших чотирьох спектральних відліків наведено у
табл.~\ref{tab:numerical}.

\begin{table}[H]
  \caption{Проміжні значення алгоритму для $x(n) = n$, рядок $r = 0$}
  \label{tab:numerical}
  \centering
  \small
  \begin{tabular}{lcccc}
    \toprule
    Матриця & Індекс 0 & Індекс 1 & Індекс 2 & Індекс 3 \\
    \midrule
    $\mathbf{A}[0][\cdot]$   & $0$       & $4$              & $8$         & $12$ \\
    $\mathbf{G}[0][\cdot]$   & $24{+}0j$ & $-8{+}8j$        & $-8{+}0j$   & $-8{-}8j$ \\
    $\mathbf{H}[0][\cdot]$   & $24{+}0j$ & $-8{+}8j$        & $-8{+}0j$   & $-8{-}8j$ \\
    $X(k)$, $k{=}0{\ldots}3$ & $120{+}0j$ & ${\approx}-8{+}40j$ & ${\approx}-8{+}19j$ & ${\approx}-8{+}12j$ \\
    \bottomrule
  \end{tabular}
\end{table}

Рядок $\mathbf{H}[0][\cdot]$ збігається з $\mathbf{G}[0][\cdot]$,
оскільки $W_{16}^{0 \cdot k_1} = 1$ для всіх $k_1$. Точне значення
$X(0) = \sum_{n=0}^{15} n = 120$ підтверджує коректність стадії~3.


% ─────────────────────────────────────────────────
\subsection{Механізми обміну даними та інтеграції}

Оскільки симулятор не має серверної частини, класичні механізми
мережевої інтеграції (REST API, WebSocket, gRPC) у ньому не
застосовуються. Замість них координація підсистем реалізується
через \emph{внутрішній подійний канал} і \emph{реактивне сховище
стану}, описані нижче.

\subsubsection*{Шина подій симуляції}

Ядро \texttt{DataflowRuntime} публікує три типи подій, що утворюють
внутрішній API симулятора (табл.~\ref{tab:events}).

\begin{table}[H]
  \caption{Події шини симуляції}
  \label{tab:events}
  \centering
  \small
  \begin{tabular}{p{2.2cm}p{5.2cm}p{5.8cm}}
    \toprule
    Подія & Умова виникнення & Підписники \\
    \midrule
    \texttt{onToken}  & Токен надійшов на дугу $e \in E$ & Heatmap черг, Metrics (queue size) \\
    \texttt{onFire}   & Активовано вузол $v \in V$ & Heatmap активності вузлів~(\ref{eq:heat_decay}) \\
    \texttt{onOutput} & Токен надійшов до \texttt{sink} & SpectrumView, Metrics (throughput, latency) \\
    \bottomrule
  \end{tabular}
\end{table}

Події реалізують шаблон Publisher/Subscriber (підрозд.~2.1):
видавець (ядро) не має прямих посилань на підписників, а лише
передає об'єкт події у список зареєстрованих обробників. Це
забезпечує ізоляцію бізнес-логіки від UI та дозволяє додавати нові
підписники (наприклад, експорт event~log у файл) без модифікації
ядра.

\subsubsection*{Реактивне сховище стану}

Стан симулятора централізований у сховищі \texttt{simStore}, що
містить такі основні групи полів:

\begin{itemize}
  \item \textit{Стан графа}: \texttt{graph: Graph},
        \texttt{queues: Map<EdgeId, Token[]>},
        \texttt{nodeBusy: Map<NodeId, number>};
  \item \textit{Стан виконання}: \texttt{status: 'IDLE' | 'RUNNING' |
        'PAUSED' | 'DONE'}, \texttt{simTime: number},
        \texttt{speed: number};
  \item \textit{Метрики}: \texttt{throughputEMA},
        \texttt{latencyEMA}, \texttt{queueSizeEMA},
        \texttt{activity: Map<NodeId, number>};
  \item \textit{Результат}: \texttt{spectrum: Complex[16]},
        \texttt{eventLog: Event[]}.
\end{itemize}

UI-компоненти підписуються на окремі поля сховища за допомогою
селекторів (pattern Selector). Компонент перемальовується лише
при зміні саме тих полів, на які він підписаний, що забезпечує
стабільні 60~fps (NF3) навіть при інтенсивному потоці подій від
ядра.

\subsubsection*{Механізм пресетів та збереження стану}

Для підтримки повторюваних демонстраційних сценаріїв (UC1 з
підрозд.~1.5) передбачено механізм пресетів. Пресет --- це JSON-об'єкт,
що містить: вхідний сигнал $\{x_n\}$, значення \texttt{latency} і
\texttt{delay}, початкову швидкість симуляції. Пресети зберігаються
у Local Storage браузера за ключем \texttt{simulator.presets} і
завантажуються через окремий компонент \texttt{PresetSelector}.
Експорт пресету у файл (формат \texttt{.json}) реалізується
стандартним Web API \texttt{Blob + URL.createObjectURL}.

Експорт event~log (UC9) реалізований аналогічно: журнал подій
серіалізується у JSON або CSV і завантажується користувачем без
передачі на сервер, що відповідає принципу конфіденційності
(підрозд.~1.4).


% ─────────────────────────────────────────────────
\subsection{Інтерфейс користувача та взаємодія із системою}

Інтерфейс симулятора спроєктовано за принципом \emph{інформаційної
ієрархії}: найбільша площа екрана відводиться під центральний
об'єкт дослідження --- граф потоку даних, --- а допоміжні панелі
(метрики, спектр, керування) розміщуються по периметру, не
перекриваючи граф.

\subsubsection*{Структура інтерфейсу}

Головний екран застосунку розділений на чотири функціональні зони:

\begin{enumerate}
  \item \textbf{Центральна зона (GraphView)} --- полотно React~Flow
        з вузлами та анімованими дугами. Користувач може перетягувати
        вузли, змінювати масштаб колесом миші, клікати по вузлу для
        перегляду деталей (поточні значення на входах/виходах).
  \item \textbf{Верхня панель (Controls)} --- кнопки «Пуск», «Пауза»,
        «Крок», «Скидання», повзунок множника швидкості ($0.1\times$
        до $10\times$), селектор пресету вхідного сигналу, поле
        редагування \texttt{latency}/\texttt{delay}.
  \item \textbf{Права панель (SpectrumView + Metrics)} --- стовпчикова
        діаграма $|X(k)|$ у лінійному або логарифмічному масштабі та
        три числових індикатори (throughput, latency, queue size).
  \item \textbf{Нижня панель (TerminalOutput)} --- журнал подій у
        реальному часі з можливістю фільтрації за типом події
        (\texttt{onToken}/\texttt{onFire}/\texttt{onOutput}) і
        експорту у файл.
\end{enumerate}

\subsubsection*{UX-логіка взаємодії}

Базовий user-flow для студента (UC3 з підрозд.~1.5) складається з
трьох кроків, доступних у межах двох кліків від головного екрана
(NF8):

\begin{enumerate}
  \item Вибір пресету вхідного сигналу (синусоїда, імпульс, шум)
        або введення власних відліків через компонент
        \texttt{InputEditor};
  \item Натискання кнопки «Пуск» --- ядро починає викликати
        \texttt{tick()} у циклі \texttt{requestAnimationFrame}
        з кроком $\mathrm{d}\tau = 16.67$~мс $\cdot$ speed;
  \item Спостереження за рухом токенів у GraphView, зміною метрик у
        Metrics та поступовим заповненням SpectrumView.
\end{enumerate}

У режимі паузи доступний покроковий режим \emph{Step}, у якому одна
кнопка викликає один виклик \texttt{tick()}: це дозволяє детально
дослідити виконання алгоритму у навчальних цілях.

\subsubsection*{Візуальні акценти на потоковій природі алгоритму}

Для підсилення сприйняття потокової природи алгоритму UI використовує
три допоміжні візуальні механізми:

\begin{itemize}
  \item \textbf{Анімація токенів.} Кожен токен відображається як
        кружечок, що переміщується по дузі в інтервалі часу
        $[t, t+\texttt{delay}]$ з плавною інтерполяцією позиції.
        На кружечку виводиться числове значення (дійсна та уявна
        частини) для прозорого спостереження обчислень.
  \item \textbf{Теплова карта активності вузлів.} Заливка кожного
        вузла модулюється значенням $a_i(\tau)$, що оновлюється за
        експоненціальним законом:
        \begin{equation}
          a_i(\tau + \mathrm{d}\tau) = \alpha \cdot a_i(\tau) + \Delta a_i,
          \quad \alpha = 0{,}92,
          \label{eq:heat_decay}
        \end{equation}
        де $\Delta a_i = 1$ при подіях \texttt{onFire} вузла $i$ та
        $\Delta a_i = 0$ в інших випадках. Вузли, що нещодавно
        часто активувалися, відображаються яскравішими («гарячими»).
  \item \textbf{Стан busy на вузлах.} Поки вузол перебуває у стані
        виконання ($v.\text{busy} > \tau$), навколо нього
        відображається кільцевий індикатор прогресу, що ілюструє
        затримку \texttt{latency} вузла.
\end{itemize}

\subsubsection*{Макет інтерфейсу}

Загальний вигляд інтерфейсу подано на рис.~\ref{fig:ui}.

\begin{figure}[H]
  \centering
  \IfFileExists{images/simulator_ui.png}{%
    \includegraphics[width=0.95\textwidth]{images/simulator_ui.png}%
  }{%
    \fbox{\parbox[c][0.24\textheight][c]{0.95\textwidth}{\centering
      Скріншот буде додано після фінального тестування.\\
      Файл: \texttt{tex/images/simulator\_ui.png}}}%
  }
  \caption{Макет інтерфейсу веб-симулятора потокового 16-точкового
           ШПФ: центральна зона GraphView, верхня панель Controls,
           права панель SpectrumView + Metrics, нижня панель
           TerminalOutput. Джерело: розроблено автором.}
  \label{fig:ui}
\end{figure}


% ─────────────────────────────────────────────────
\subsection*{Висновки до розділу~2}

У другому розділі спроєктовано програмне рішення, що відповідає
вимогам, сформульованим у підрозд.~1.5.

\emph{Обрано архітектурний підхід} клієнтського SPA без серверної
частини, що забезпечує одночасне виконання вимог NF2 (відгук
$\le 100$~мс), NF4 (портативність Chrome/Firefox/Edge), NF7 (робота
без встановлення) і NF5 (детермінованість). Вибір реактивного
рендерингу та централізованого сховища стану обґрунтований як
компроміс між продуктивністю та супроводжуваністю коду. Застосовано
шаблони проєктування Observer, State, Strategy, Factory і Facade;
зв'язок між вимогами та рішеннями показано у
табл.~\ref{tab:req_to_tech}.

\emph{Побудовано математичну модель} алгоритму ШПФ radix-4 для
$N = 16$ у вигляді тристадійної $4 \times 4$ декомпозиції
(\ref{eq:stage1})--(\ref{eq:stage3}), що потребує 9~нетривіальних
комплексних множень і 80~комплексних додавань проти 256~множень
прямого ДПФ. Сформульовано формальну модель потоку даних із
правилом активації та часовою семантикою, подано псевдокод
планувальника (рис.~\ref{fig:tick_pseudo}).

\emph{Спроєктовано архітектуру системи} на рівнях C4 Context і
Container з декомпозицією на чотири підсистеми (Simulation Core,
Graph Generator, State Store, UI Layer). Деталізовано граф
обчислювача: 56~вузлів чотирьох типів і 64~дуги з однаковою
затримкою 600~одиниць, критичним шляхом $T_{\mathrm{kr}} = 2200$
і коефіцієнтом паралелізму $K_p \approx 2{,}9$.

\emph{Визначено структури даних:} \texttt{Token}, \texttt{Edge},
\texttt{NodeSpec}, \texttt{Graph}, що є програмними відповідниками
формальних об'єктів DFG. Комплексні числа подаються у форматі
IEEE~754 double, що гарантує запас точності відносно
$\varepsilon = 10^{-9}$ (NF1).

\emph{Обрано механізми обміну даними:} внутрішня шина подій
(\texttt{onToken}, \texttt{onFire}, \texttt{onOutput}) для ізоляції
UI від ядра, реактивне сховище стану та пресети у Local~Storage
для повторюваних сценаріїв демонстрації.

\emph{Спроєктовано інтерфейс користувача} з інформаційною ієрархією
(центральний GraphView, периферійні Controls, SpectrumView, Metrics,
TerminalOutput), трьома візуальними механізмами для підсилення
потокової природи алгоритму (анімація токенів, теплова карта
активності, індикатор busy) та базовим user-flow у межах двох
кліків (NF8).

Описані проєктні рішення утворюють повне технічне завдання для
реалізації, що подається у розділі~3.
\newpage

% --- Розділ 3 ---
% chapters/chapter3.tex
\section{РЕАЛІЗАЦІЯ ПРОГРАМНОГО РІШЕННЯ}

% Розділ 3 згідно нових методичних вказівок (15-20 стор.):
% Для типу роботи «ПЗ / Веб / ІС» підрозділи 3.4 «Дані та їх підготовка»
% (обов'язковий лише для ML/ШІ) та 3.5 «Апаратна частина» (обов'язковий
% лише для IoT / embedded) не застосовуються. Використано послідовну
% нумерацію без пропусків:
%
% 3.1 Опис технологій та інструментів розробки
% 3.2 Структура проєкту та організація вихідного коду
% 3.3 Реалізація ключових модулів та алгоритмів
% 3.4 Обробка виняткових ситуацій та забезпечення надійності
% 3.5 Розгортання системи
% 3.6 Тестування програмного забезпечення
% 3.7 Висновки до розділу 3

% -------------------------------------------------
\subsection{Опис технологій та інструментів розробки}

У підрозд.~2.1 обґрунтовано архітектурні рішення: клієнтський SPA,
реактивний рендеринг, централізоване сховище стану, спеціалізована
бібліотека графової візуалізації. У цьому підрозділі наведено
конкретні технології, бібліотеки, їхні версії та середовище розробки,
що реалізують зазначені архітектурні рішення.

\subsubsection*{Мови програмування та виконуване середовище}

Основною мовою реалізації обрано \textbf{TypeScript~5.x} --- надбудову
над JavaScript зі статичною типізацією. Мова компілюється у
стандартизований \textbf{ECMAScript~2020 (ES11)}, який виконується
усіма сучасними браузерами без додаткових polyfill'ів, що
задовольняє вимогу NF4. Статична типізація дозволяє виявляти на
етапі компіляції помилки передачі об'єкта \texttt{Token} замість
числа чи неправильної форми комплексного значення, що критично для
дотримання NF1 (точність $\varepsilon \le 10^{-9}$).

Цільовим середовищем виконання є браузерний рушій JavaScript
(V8 у Chrome/Edge, SpiderMonkey у Firefox) з підтримкою
\texttt{requestAnimationFrame}, Web Workers, Canvas API та Local
Storage.

\subsubsection*{Фреймворки та бібліотеки}

Повний перелік використаних бібліотек із версіями та призначенням
наведений у табл.~\ref{tab:tech_stack}.

\begin{table}[H]
  \caption{Технологічний стек реалізації}
  \label{tab:tech_stack}
  \centering
  \small
  \begin{tabular}{p{3.2cm}p{1.6cm}p{8.2cm}}
    \toprule
    Технологія & Версія & Призначення \\
    \midrule
    TypeScript & 5.x & Статично типізована мова реалізації \\
    Next.js & 14.x & React-фреймворк, статичний експорт SPA \\
    React & 18.x & Реактивний рендеринг UI-компонентів \\
    React Flow (\texttt{@xyflow/react}) & 12.x & Інтерактивна графова візуалізація DFG \\
    Zustand & 4.x & Централізоване реактивне сховище стану \\
    Comlink & 4.x & RPC-обгортка для Web Workers \\
    fft.js & 4.x & Референтна реалізація ДПФ для верифікації \\
    Jest & 29.x & Фреймворк модульного тестування \\
    ESLint & 8.x & Статичний аналіз коду, правила стилю \\
    Prettier & 3.x & Автоматичне форматування коду \\
    \bottomrule
  \end{tabular}
\end{table}

\textbf{Next.js} забезпечує маршрутизацію, оптимізоване завантаження
ресурсів і статичну збірку без серверної частини (режим
\texttt{output: 'export'}). Усі сторінки симулятора позначені
директивою \texttt{'use client'}, оскільки синхронне клієнтське
виконання \texttt{requestAnimationFrame}-циклу несумісне зі
server-side rendering.

\textbf{React Flow} надає готові механізми перетягування вузлів,
масштабування, анімованих ребер і кастомних типів вузлів. Його
абстрактна модель \texttt{Node}/\texttt{Edge} безпосередньо
відображається на формальну модель DFG (підрозд.~2.4).

\textbf{Zustand} --- легка альтернатива Redux без шаблонного коду
редьюсерів. Забезпечує реактивне сховище стану, до якого компоненти
підписуються через хук \texttt{useSimStore(selector)} з мемоїзацією.

\textbf{fft.js} --- компактна JavaScript-реалізація ДПФ, що
використовується виключно як референтна для верифікації
(підрозд.~3.6) і не залучається до шляху виконання основного
алгоритму, що забезпечує незалежність результату симулятора від
зовнішньої бібліотеки.

\subsubsection*{Середовище розробки та інструментарій}

Розробка виконується у редакторі \textbf{Visual Studio Code}
з розширеннями ESLint, Prettier і TypeScript. Pre-commit-хуки
реалізовано через \texttt{husky} + \texttt{lint-staged}: перед
кожним коммітом автоматично запускаються ESLint та Prettier
для змінених файлів.

Менеджер пакетів --- \textbf{pnpm}, що забезпечує швидке встановлення
залежностей і ефективне використання дискового простору через
hard-linking. Файлом залежностей є \texttt{package.json}; файл
\texttt{pnpm-lock.yaml} гарантує детермінованість збірки на різних
машинах.

\subsubsection*{Система керування версіями та репозиторій}

Використовується \textbf{Git} зі стратегією \textit{trunk-based
development}: основна розробка ведеться у гілці \texttt{main}, для
великих змін створюються короткоживучі feature-гілки. Репозиторій
розміщений на GitHub, що забезпечує доступність для науково-педагогічного
комітету БКР.

Опційно налаштовано \textbf{GitHub Actions} для автоматичного запуску
модульних тестів і лінтерів на кожному push до \texttt{main}
(див. підрозд.~3.5).


% -------------------------------------------------
\subsection{Структура проєкту та організація вихідного коду}

\subsubsection*{Верхньорівневе дерево репозиторію}

Структура кореневого каталогу проєкту організована за стандартом
Next.js з додатковими директоріями для тестів і документації
(рис.~\ref{fig:project_tree}).

\begin{figure}[H]
\begin{verbatim}
simulator/
+-- public/                  # статичні ресурси (SVG-іконки, favicons)
+-- src/
|   +-- app/                 # Next.js App Router
|   |   +-- layout.tsx
|   |   \-- page.tsx
|   +-- core/                # ядро симуляції (незалежне від UI)
|   |   +-- types.ts
|   |   +-- scheduler.ts
|   |   \-- buffers.ts
|   +-- graph/               # генерація та візуалізація графа
|   |   +-- generateDFT4x4.ts
|   |   +-- layout4x4.ts
|   |   +-- GraphView.tsx
|   |   \-- TokenEdge.tsx
|   +-- components/          # UI-компоненти
|   |   +-- SpectrumView.tsx
|   |   +-- Metrics.tsx
|   |   +-- Controls.tsx
|   |   +-- TerminalOutput.tsx
|   |   \-- nodes/
|   |       +-- SourceNode.tsx
|   |       +-- SinkNode.tsx
|   |       +-- DFT4Node.tsx
|   |       \-- TwiddleNode.tsx
|   +-- store/               # централізоване сховище стану
|   |   \-- simStore.ts
|   +-- utils/               # допоміжні утиліти
|   |   +-- math.ts
|   |   \-- compare.ts
|   \-- workers/             # Web Workers
|       \-- fft.worker.ts
+-- tests/                   # Jest-тести
|   +-- math.test.ts
|   +-- scheduler.test.ts
|   \-- compare.test.ts
+-- package.json
+-- tsconfig.json
+-- next.config.js
\-- .eslintrc.js
\end{verbatim}
\caption{Структура директорій проєкту. Джерело: розроблено автором.}
\label{fig:project_tree}
\end{figure}

\subsubsection*{Принцип декомпозиції}

Декомпозиція на модулі виконана за принципом \emph{розподілу за
шарами відповідальності}, що відповідає архітектурі Container з
підрозд.~2.3:

\begin{itemize}
  \item \texttt{src/core/} --- ядро симуляції без будь-яких
        залежностей від UI; може бути імпортоване у Node.js для
        batch-обчислень і використовується в unit-тестах;
  \item \texttt{src/graph/} --- генерація структури DFG і компоненти
        інтерактивної візуалізації на React Flow;
  \item \texttt{src/components/} --- UI-компоненти, що не залежать
        від React Flow (панелі керування, метрики, спектр);
  \item \texttt{src/store/} --- реактивне сховище стану та
        інтеграція з ядром симуляції;
  \item \texttt{src/utils/} --- чисті функції (математика, верифікація)
        без побічних ефектів;
  \item \texttt{src/workers/} --- Web Worker для фонових обчислень.
\end{itemize}

Така організація забезпечує можливість ізольованого тестування ядра
симуляції без браузерного середовища та спрощує пошук місць
внесення змін при підтримці коду.

\subsubsection*{Відповідність компонентів архітектури та модулів коду}

Таблиця~\ref{tab:arch_to_code} формалізує зв'язок між компонентами
архітектури з підрозд.~2.3 та модулями вихідного коду.

\begin{table}[H]
  \caption{Відповідність «Компонент $\to$ Модуль/файл»}
  \label{tab:arch_to_code}
  \centering
  \small
  \begin{tabular}{p{3.5cm}p{5.5cm}p{4.2cm}}
    \toprule
    Компонент архітектури (2.3) & Модуль / файл & Роль \\
    \midrule
    Simulation Core & \texttt{src/core/scheduler.ts} & Клас \texttt{DataflowRuntime}, метод \texttt{tick()} \\
    Simulation Core & \texttt{src/core/buffers.ts} & FIFO-черги на дугах \\
    Simulation Core & \texttt{src/core/types.ts} & Типи \texttt{Token}, \texttt{Edge}, \texttt{NodeSpec} \\
    Graph Generator & \texttt{src/graph/generateDFT4x4.ts} & Побудова DFG \\
    Graph Generator & \texttt{src/graph/layout4x4.ts} & Координати вузлів \\
    State Store & \texttt{src/store/simStore.ts} & Zustand-сховище \\
    UI Layer (GraphView) & \texttt{src/graph/GraphView.tsx} & React Flow + rAF-цикл \\
    UI Layer (SpectrumView) & \texttt{src/components/SpectrumView.tsx} & Canvas-діаграма $|X(k)|$ \\
    UI Layer (Metrics) & \texttt{src/components/Metrics.tsx} & EMA-метрики \\
    UI Layer (Controls) & \texttt{src/components/Controls.tsx} & Пуск/пауза/швидкість \\
    UI Layer (вузли) & \texttt{src/components/nodes/*.tsx} & SourceNode, SinkNode, DFT4Node, TwiddleNode \\
    Math \& verify & \texttt{src/utils/math.ts}, \texttt{compare.ts} & Комплексна арифметика, верифікація \\
    \bottomrule
  \end{tabular}
\end{table}

\subsubsection*{Ключові конфігураційні файли}

\texttt{tsconfig.json} визначає налаштування компілятора TypeScript:
увімкнені \texttt{strict}, \texttt{noImplicitAny},
\texttt{strictNullChecks} --- для максимальної безпеки типів.
Цільовий стандарт --- ES2020.

\texttt{next.config.js} встановлює \texttt{output: 'export'} для
статичної збірки та налаштовує базовий шлях. Функцій і API немає,
усі сторінки клієнтські.

\texttt{.eslintrc.js} налаштовує правила ESLint з розширеннями
\texttt{@typescript-eslint}, \texttt{react-hooks},
\texttt{@next/next}. Заборонено \texttt{any}, \texttt{console.log}
у production, не імпортовані змінні, побічні ефекти у хуках
без залежностей.

\subsubsection*{Стандарти кодування}

Код документується у стилі \textbf{JSDoc}, з особливим акцентом
на публічних методах \texttt{DataflowRuntime}, \texttt{simStore} і
експортованих утилітах. Типи та інтерфейси описуються безпосередньо
у \texttt{types.ts}. Ідентифікатори --- англомовні, камел-кейс для
змінних і функцій, паскаль-кейс для типів і React-компонентів.
Максимальна довжина рядка --- 100~символів, ця межа автоматично
контролюється Prettier.


% -------------------------------------------------
\subsection{Реалізація ключових модулів та алгоритмів}

\subsubsection*{Клас DataflowRuntime}

Клас \texttt{DataflowRuntime} (модуль \texttt{src/core/\allowbreak scheduler.ts})
є програмною реалізацією абстрактної машини потоку даних
з підрозд.~2.2~\cite{veen1986,dennis1975}. Він відповідає за
виконання правила активації вузлів та управління FIFO-чергами
на дугах графа.

Клас зберігає такі внутрішні структури даних:
\begin{itemize}
  \item незмінне посилання на граф \texttt{Graph} (вузли і дуги);
  \item відображення \texttt{Map<edgeId, FIFO<Token>>} --- окрема
        черга токенів для кожної дуги;
  \item \texttt{readyAt: Map<nodeId, number>} --- час, раніше якого
        вузол не може активуватись повторно (моделює \texttt{latency});
  \item \texttt{inbox: Map<nodeId, Map<portId, Token>>} --- токени,
        що вже доставлені до вузла, але ще не поглинуті;
  \item обробники подій: \texttt{onToken}, \texttt{onFire},
        \texttt{onBeforeFire}, \texttt{onOutput}.
\end{itemize}

Логічний час надається ззовні через функцію
\texttt{getTime: () => number}, що передається у конструктор. У
симуляторі цією функцією є \texttt{() => store.simTime}, завдяки чому
годинник ядра повністю відокремлений від астрономічного часу.

\subsubsection*{Метод tick() --- алгоритм одного кроку симуляції}

Метод \texttt{tick(dt, maxFires)} виконує єдиний крок потокового
обчислення і реалізує правило активації (firing rule), описане у
підрозд.~2.2~\cite{veen1986,dennis1975}. Алгоритм складається
з чотирьох послідовних кроків:

\begin{enumerate}
  \item \textbf{Доставка токенів.} Для кожної дуги $e$ перевіряється,
    чи найстарший токен у черзі $\mathrm{buf}[e.\mathrm{id}]$
    задовольняє умові доставки:
    \begin{equation}
      t_{\text{now}} - \mathrm{token}.t \;\geq\; e.\mathrm{delay},
      \label{eq:delivery}
    \end{equation}
    де $t_{\text{now}} = \texttt{simTime}$. Якщо умова~(\ref{eq:delivery})
    виконана, токен видаляється з черги і поміщається у \texttt{inbox}
    відповідного вузла-приймача; емітується подія \texttt{onToken}.

  \item \textbf{Перевірка готовності вузла.} Вузол $n$ є
    \emph{готовим} до активації, якщо одночасно виконуються дві умови:
    \begin{itemize}
      \item[(а)] у \texttt{inbox}$[n]$ є токен на \emph{кожному}
        вхідному порті: $\forall p \in n.\mathrm{inPorts}:
        \texttt{inbox}[n].\mathrm{has}(p)$;
      \item[(б)] вузол вільний від блокування:
        $\texttt{readyAt}[n] \leq t_{\text{now}}$.
    \end{itemize}

  \item \textbf{Активація вузла.} Для кожного готового вузла
    (не більше \texttt{maxFires} за один виклик):
    \begin{itemize}
      \item[3.1] токени поглинаються з \texttt{inbox}$[n]$ і
        передаються у метод \texttt{execute(node, inputs)};
      \item[3.2] виконується логіка вузла за типом
        \texttt{NodeKind}, обчислюються вихідні токени;
      \item[3.3] час блокування оновлюється:
        $\texttt{readyAt}[n] \leftarrow t_{\text{now}} +
        n.\mathrm{latency}$;
      \item[3.4] емітуються події \texttt{onBeforeFire} і
        \texttt{onFire}.
    \end{itemize}

  \item \textbf{Розміщення вихідних токенів.} Для кожного вихідного
    токена знаходяться всі дуги з відповідним портом-джерелом; токен
    поміщається у чергу кожної такої дуги з часом народження:
    \begin{equation}
      \mathrm{token}.t = t_{\text{now}} + n.\mathrm{latency},
      \label{eq:emit_time}
    \end{equation}
    що гарантує правильну затримку доставки на наступному кроці.
\end{enumerate}

Такий порядок кроків забезпечує детермінованість: при фіксованому
порядку обходу вузлів і дуг результат симуляції не залежить від
реального часу виконання у браузері (NF5).

\subsubsection*{FIFO-черга (buffers.ts)}

Клас \texttt{FIFO<T>} (модуль \texttt{src/core/buffers.ts}) ---
узагальнена черга на базі масиву з методами \texttt{push(item)},
\texttt{pop(): T | undefined}, \texttt{peek(): T | undefined} та
властивістю \texttt{size: number}. Одна черга ставиться у
відповідність кожній дузі графа; черги ініціалізуються у конструкторі
\texttt{DataflowRuntime} ітерацією по \texttt{graph.edges}.

Асимптотична складність операцій: \texttt{push} --- $O(1)$
амортизований, \texttt{pop} --- $O(n)$ у найгіршому випадку (через
зсув масиву). При характерних для симулятора розмірах черг
($n \le 16$) це не становить проблеми; для масштабування на більші
$N$ заплановано перехід на циклічний буфер з $O(1)$-операціями.

\subsubsection*{Логічний час і speed multiplier}

Логічний час \texttt{simTime} збільшується методом
\texttt{tickSimTime()} пропорційно до астрономічного часу та
параметра \texttt{speed}:

\begin{equation}
  \texttt{simTime}_{i+1} = \texttt{simTime}_{i} +
    \Delta t_{\text{wall}} \cdot \texttt{speed},
  \label{eq:simtime}
\end{equation}

де $\Delta t_{\text{wall}}$ --- реальний проміжок між кадрами
(у мілісекундах), а \texttt{speed} --- множник у діапазоні
$[0{,}1;\; 10]$. Відокремлення логічного часу від астрономічного
дозволяє досліджувати поведінку алгоритму у довільному
масштабі без зміни логіки активації.

\subsubsection*{Реалізація вузлів за типом}

Метод \texttt{execute(node, inputs)} реалізує шаблон Strategy
(підрозд.~2.1) через перемикач за \texttt{node.kind}.

\textbf{SOURCE.} Вузол-джерело не має вхідних портів; його токени
вводяться ззовні через виклик
\texttt{runtime.emitFrom(nodeId, 'out', token)}. Токен містить
значення $x(n) = \{\,re, im\,\}$ і поля
$\texttt{t} = \texttt{originT} = \texttt{simTime}$.

\textbf{DFT4.} Вузол приймає 4~токени через порти \texttt{in0}
\ldots \texttt{in3} і застосовує алгоритм метелика
radix-4~(\ref{eq:dft4_butterfly}) з підрозд.~2.2. Лістинг~\ref{lst:dft4}
наводить реалізацію:

\begin{figure}[H]
\begin{verbatim}
function computeDFT4(x: Complex[4]): Complex[4] {
  const E0 = cplxAdd(x[0], x[2]);   // x0 + x2
  const E1 = cplxSub(x[0], x[2]);   // x0 - x2
  const O0 = cplxAdd(x[1], x[3]);   // x1 + x3
  const O1 = cplxSub(x[1], x[3]);   // x1 - x3
  return [
    cplxAdd(E0, O0),                // X0 = E0 + O0
    cplxSub(E1, cplxMulJ(O1)),      // X1 = E1 - j*O1
    cplxSub(E0, O0),                // X2 = E0 - O0
    cplxAdd(E1, cplxMulJ(O1))       // X3 = E1 + j*O1
  ];
}
\end{verbatim}
\caption{Реалізація 4-точкового ДПФ-метелика без нетривіальних
         множень. Функція \texttt{cplxMulJ(z) = \{-z.im, z.re\}}
         реалізує множення на $j$ через перестановку. Джерело:
         розроблено автором.}
\label{lst:dft4}
\end{figure}

Вузол породжує 4~вихідних токени через порти
\texttt{out0} \ldots \texttt{out3}. Жодного нетривіального множення
не виконується.

\textbf{TWIDDLE.} Вузол приймає 1~токен і виконує комплексне
множення на поворотний множник $W_{16}^k = e^{-j 2\pi k/16}$, де
ціле число $k$ зберігається у \texttt{NodeSpec.params.twiddle.k}.
Формули обчислення:

\begin{equation}
  \begin{aligned}
    \mathrm{re}' &= z.\mathrm{re} \cos\theta - z.\mathrm{im} \sin\theta, \\
    \mathrm{im}' &= z.\mathrm{re} \sin\theta + z.\mathrm{im} \cos\theta,
  \end{aligned}
  \quad \theta = -\frac{2\pi k}{16}.
  \label{eq:twiddle_impl}
\end{equation}

При $k = 0$ маємо $W_{16}^0 = 1$, і множення є тривіальним ---
вхідний токен передається на вихід без змін.

\textbf{SINK.} Вузол приймає 1~токен, передає його значення у
\texttt{store.sinks[nodeId]} через серіалізацію
\texttt{JSON.stringify} і емітує подію \texttt{onOutput}.
Одночасно обчислюється наскрізна затримка:

\begin{equation}
  \mathrm{latency} = \texttt{simTime} - \mathrm{token}.\texttt{originT},
  \label{eq:lat_calc}
\end{equation}

і оновлюється EMA-метрика \texttt{latencyEmaMs}.

\subsubsection*{Функція generateDFT4x4()}

Функція \texttt{generateDFT4x4(latency): Graph} у модулі
\texttt{src/graph/\allowbreak generateDFT4x4.ts} реалізує шаблон Factory
(підрозд.~2.1) і будує DFG за алгоритмом:

\begin{enumerate}
  \item створюються 16~вузлів \texttt{source}
        (\texttt{src0} \ldots \texttt{src15}), кожен з одним
        вихідним портом \texttt{out};
  \item створюються 4~вузли \texttt{dft4} першого ярусу
        (\texttt{row-0} \ldots \texttt{row-3}) з 4~входами та
        4~виходами і \texttt{latency} $= 400$;
  \item додаються 16~дуг \texttt{src}$_i \to$
        \texttt{row-}$r$\texttt{.in}$s$, де $i = r + 4s$,
        \texttt{delay} $= 600$;
  \item створюються 16~вузлів \texttt{twiddle}
        (\texttt{tw-}$r$\texttt{-}$k_1$) з параметром
        $\{N: 16,\; k: r \cdot k_1\}$ і \texttt{latency} $= 200$;
  \item додаються 16~дуг \texttt{row-}$r$\texttt{.out}$k_1 \to$
        \texttt{tw-}$r$\texttt{-}$k_1$\texttt{.in};
  \item створюються 4~вузли \texttt{dft4} другого ярусу
        (\texttt{col-0} \ldots \texttt{col-3});
  \item додаються 16~дуг \texttt{tw-}$r$\texttt{-}$k_1$\texttt{.out}
        $\to$ \texttt{col-}$k_1$\texttt{.in}$r$;
  \item створюються 16~вузлів \texttt{sink}
        (\texttt{snk0} \ldots \texttt{snk15});
  \item додаються 16~дуг \texttt{col-}$k_1$\texttt{.out}$k_2 \to$
        \texttt{snk}$_{k_2 + 4 k_1}$\texttt{.in}.
\end{enumerate}

Повертається об'єкт \texttt{Graph} із 56~вузлами і 64~дугами. Функція
\texttt{layoutDFT4x4()} (модуль \texttt{layout4x4.ts}) обчислює
координати вузлів для React Flow: граф розподіляється на п'ять
вертикальних ярусів з кроком \texttt{stageGap} $= 280$ пікселів;
вузли в межах ярусу рівномірно розподіляються по осі~$y$ з кроком
\texttt{rowGap} $= 140$ пікселів.

\subsubsection*{Сховище стану simStore (Zustand)}

Сховище \texttt{simStore} реалізовано функцією \texttt{create()}
з Zustand і містить поля стану та методи-команди
(табл.~\ref{tab:store}).

\begin{table}[H]
  \caption{Ключові поля та методи Zustand-сховища}
  \label{tab:store}
  \centering
  \small
  \begin{tabular}{p{3.2cm}p{3.2cm}p{5.8cm}}
    \toprule
    Поле / метод & Тип & Призначення \\
    \midrule
    \texttt{graph} & \texttt{Graph} & Поточна топологія DFG \\
    \texttt{runtime} & \texttt{DataflowRuntime\textbar null} & Ядро симуляції \\
    \texttt{running} & \texttt{boolean} & Стан rAF-циклу \\
    \texttt{speed} & \texttt{number} & Множник швидкості $[0{,}1; 10]$ \\
    \texttt{simTime} & \texttt{number} & Поточний логічний час \\
    \texttt{tokensByEdge} & \texttt{Record<id, Token[]>} & Токени для анімації \\
    \texttt{sinks} & \texttt{Record<id, Complex>} & Спектральні відліки $X(k)$ \\
    \texttt{nodeActivity} & \texttt{Record<id, number>} & Теплова карта $a_i$ \\
    \texttt{metrics} & \texttt{MetricsState} & EMA-метрики \\
    \texttt{mismatches} & \texttt{Record<id, boolean>} & Результат верифікації \\
    \texttt{setGraph(g)} & \texttt{(Graph) => void} & Ініціалізація нового графа і \texttt{runtime} \\
    \texttt{injectData(v)} & \texttt{(Complex[16]) => void} & Введення вхідного вектора \\
    \texttt{start()/stop()/step()} & \texttt{() => void} & Керування rAF-циклом \\
    \texttt{applyPreset(k)} & \texttt{(PresetKind) => void} & Вбудовані пресети вхідних сигналів \\
    \bottomrule
  \end{tabular}
\end{table}

Компоненти UI підписуються на окремі зрізи стану через хук
\texttt{useSimStore(selector)} з мемоїзацією: наприклад, компонент
\texttt{Metrics} підписується лише на поле \texttt{metrics} і
перемальовується тільки при його зміні, що забезпечує стабільні
60~fps (NF3).

\subsubsection*{Анімація токенів (TokenEdge)}

Компонент \texttt{TokenEdge} (модуль \texttt{src/graph/\allowbreak TokenEdge.tsx})
відображає кожну дугу як анімовану стрілку з кружечками-токенами.
Для кожного токена у \texttt{store.tokensByEdge[edgeId]} обчислюється
відносний прогрес руху:

\begin{equation}
  \mathrm{progress} = \frac{\texttt{simTime} - \mathrm{token}.t_0}
                           {\mathrm{token}.\texttt{delay}}
                    \in [0, 1],
  \label{eq:progress}
\end{equation}

де $t_0$ --- момент виникнення токена на дузі. Поточна координата
визначається лінійною інтерполяцією між центрами вузлів-кінців:

\begin{equation}
  \mathbf{r}(\mathrm{progress}) =
    (1 - \mathrm{progress}) \cdot \mathbf{r}_{\text{src}}
    + \mathrm{progress} \cdot \mathbf{r}_{\text{tgt}}.
  \label{eq:lerp}
\end{equation}

Колір кружечка кодує тип вузла-джерела, значення $\{re, im\}$
виводиться всередині кружечка для прозорого спостереження.

\subsubsection*{Теплова карта активності вузлів}

Теплова карта реалізується у \texttt{simStore} через механізм
експоненціального затухання~(\ref{eq:heat_decay}) з підрозд.~2.6.
Функція \texttt{bumpActivity(nodeId)} викликається з обробника
події \texttt{onFire} та інкрементує лічильник відповідного вузла.
У кожному кадрі анімаційного циклу викликається
\texttt{decayActivity()}, що множить усі лічильники на
$\alpha = 0{,}92$. Поточне значення $a_i$ відображається як
яскравість заливки відповідного вузла через CSS-клас, що
обирається у модулі \texttt{\_shared.ts}.

\subsubsection*{Асинхронні обчислення у Web Worker}

Для уникнення блокування головного потоку при масштабуванні на
більші $N$ реалізовано Web Worker у файлі
\texttt{src/workers/fft.worker.ts}. Через бібліотеку Comlink
функція \texttt{applyTwiddle(tokens, factors)} приймає масив
токенів і масив поворотних множників, виконує пакетне комплексне
множення у фоновому потоці та повертає результат. У базовій
конфігурації $N = 16$ Web Worker не задіяний; він зарезервований
для подальшого розширення.

\subsubsection*{Цикл requestAnimationFrame у GraphView}

Головний анімаційний цикл запускається у хуку \texttt{useEffect}
компонента \texttt{GraphView} після ініціалізації \texttt{runtime}.
На кожному кадрі:

\begin{enumerate}
  \item якщо \texttt{store.running = true}, викликається
        \texttt{tickSimTime()}, потім
        \texttt{runtime.tick(dt, maxFires)}; у режимі
        \texttt{'single-fire'} \texttt{maxFires} $= 1$, у режимі
        \texttt{'run'} --- \texttt{Infinity};
  \item викликаються \texttt{sampleQueues()} та
        \texttt{decayActivity()} для оновлення метрик і теплової
        карти;
  \item щосекунди агрегуються лічильники \texttt{firesInWindow} для
        обчислення пропускної здатності;
  \item цикл перезапускається через \texttt{requestAnimationFrame}.
\end{enumerate}

Такий дизайн гарантує, що ядро симуляції виконує один детермінований
крок на один кадр браузера, а UI оновлюється з частотою 60~Гц.


% -------------------------------------------------
\subsection{Обробка виняткових ситуацій та забезпечення надійності}

Оскільки симулятор не взаємодіє з мережею, серверами чи файловою
системою, класичні джерела збоїв (мережеві тайм-аути, транзакційні
відмови, I/O-помилки) для нього не характерні. Водночас існують
специфічні категорії помилок, пов'язаних з валідацією вхідних даних,
цілісністю графа симуляції та обмеженнями браузерного середовища.

\subsubsection*{Класифікація типових помилок}

Помилки, що можуть виникнути у симуляторі, зведено у
табл.~\ref{tab:errors}.

\begin{table}[H]
  \caption{Класифікація помилок симулятора та підходи до їх обробки}
  \label{tab:errors}
  \centering
  \small
  \begin{tabular}{p{2.8cm}p{4.0cm}p{5.6cm}}
    \toprule
    Категорія & Приклади & Обробка \\
    \midrule
    Валідаційні & Введення не-числового значення у \texttt{SourceNode}, множник швидкості поза $[0{,}1; 10]$ & Clamp до допустимого діапазону, підсвічування червоним \\
    Транзієнтні & Пропуск кадру при перевантаженні CPU & \texttt{maxFires} обмежує кількість активацій за один \texttt{tick()} \\
    Структурні & Відсутність вхідного токена на порті вузла & Вузол не активується (правило готовності); статичний аналіз графа виключає deadlock \\
    Чисельні & Втрата точності при $|X(k)| > 10^{12}$ & Використання IEEE 754 double; попередження при виході за рамки \\
    Браузерні & Відсутність Web Worker API & Graceful degradation: обчислення у головному потоці \\
    \bottomrule
  \end{tabular}
\end{table}

\subsubsection*{Валідація вхідних даних}

Компонент \texttt{SourceNode} використовує повзунки з фіксованими
межами $[-10; 10]$ для \texttt{re} та \texttt{im}, що виключає
введення нечислових значень. Компонент \texttt{Controls} аналогічно
обмежує множник \texttt{speed} діапазоном $[0{,}1; 10]$ за допомогою
\texttt{Math.max(0.1, Math.min(10, value))}. Параметри затримок
\texttt{latency} і \texttt{delay} валідуються на додатність і
цілочисельність; у разі порушення викидається виняток
\texttt{InvalidGraphError} з описовою причиною.

\subsubsection*{Відсутність ризику deadlock у DFG}

Оскільки граф симулятора є статичним орієнтованим ациклічним
графом (DAG) з детермінованою SDF-семантикою (підрозд.~1.2), ризик
взаємного блокування принципово відсутній: відсутність циклів
гарантує, що готовність вузлів поступово поширюється від
\texttt{source} до \texttt{sink}. Ця властивість перевіряється на
етапі побудови графа функцією \texttt{validateDAG(graph)} через
топологічне сортування методом Кана.

\subsubsection*{Логування та трасування}

Усі значущі події ядра (\texttt{onToken}, \texttt{onFire},
\texttt{onOutput}) акумулюються у \texttt{store.eventLog} із
часовими мітками \texttt{simTime} та \texttt{performance.now()}.
Компонент \texttt{TerminalOutput} відображає журнал у реальному
часі з можливістю фільтрації за типом події. Експорт журналу у
файл (CSV/JSON) реалізований через Web API
\texttt{Blob + URL.createObjectURL} без передачі на сервер.

Рівні логування (\texttt{debug}, \texttt{info}, \texttt{warn},
\texttt{error}) налаштовуються через прапорець
\texttt{store.logLevel}; у production-збірці рівень
\texttt{debug} вимкнений за замовчуванням.

\subsubsection*{Graceful degradation та оптимізація}

\textbf{Fallback при відсутності Web Worker.} Якщо браузер не
підтримує Web Workers, функція \texttt{applyTwiddle()} виконується
синхронно у головному потоці; для $N = 16$ це не впливає на
частоту кадрів.

\textbf{Оптимізація перерендерингу.} Селектори Zustand
мемоїзовані функцією \texttt{shallow}, що запобігає зайвим
перемалюванням компонентів при зміні несуттєвих полів сховища.
Важкі розрахунки (наприклад, обчислення координат токенів на
ребрах) виконуються лише для видимої області через перевірку
\texttt{ReactFlow.viewportBounds}.

\textbf{Профілювання.} Для розробника у режимі debug доступна панель
з часовими показниками: тривалість \texttt{tick()}, тривалість
рендеру React, інтервали між кадрами. Середня тривалість одного
кроку \texttt{tick()} для $N = 16$ становить $< 1$~мс, що залишає
понад 15~мс запасу у бюджеті кадру (16.67~мс при 60~fps).


% -------------------------------------------------
\subsection{Розгортання системи}

Розгортання симулятора спрощене порівняно з класичними
веб-застосунками завдяки відсутності серверної частини та бази
даних. Зібраний статичний бандл може бути розміщений на будь-якому
файловому хостингу, що віддає статичні ресурси.

\subsubsection*{Середовище розгортання}

Цільовим середовищем є \textbf{статичний веб-хостинг}: GitHub Pages,
Vercel (статичний режим), Netlify або локальний HTTP-сервер
розробника. Серверна інфраструктура, бази даних, сесійне сховище
та автентифікація не застосовуються.

\subsubsection*{Процедура збірки}

Збірка виконується стандартною командою Next.js:

\begin{figure}[H]
\begin{verbatim}
$ pnpm install            # встановлення залежностей
$ pnpm run build          # next build -> .next/
$ pnpm run export         # next export -> out/
$ pnpm run lint           # ESLint
$ pnpm test               # Jest-тести
\end{verbatim}
\caption{Команди збірки та перевірки проєкту. Джерело: розроблено
         автором.}
\label{lst:build}
\end{figure}

У результаті у директорії \texttt{out/} формується статичний
бандл: HTML-сторінка, JavaScript-бандл із коду застосунку (через
Turbopack/webpack із оптимізаціями tree-shaking і minification) та
CSS-файли. Розмір gzip-бандла становить орієнтовно 250~КБ.

\subsubsection*{CI/CD-пайплайн (опційно)}

Для автоматизації перевірок на кожному push до основної гілки
налаштовано GitHub Actions із workflow \texttt{.github/workflows/ci.yml}:

\begin{itemize}
  \item job \texttt{lint} --- запуск ESLint і TypeScript-компіляції
        у режимі \texttt{--noEmit} (перевірка типів без генерації
        файлів);
  \item job \texttt{test} --- запуск Jest-тестів з генерацією звіту
        про покриття;
  \item job \texttt{build} --- перевірка, що збірка проходить без
        помилок.
\end{itemize}

Стратегія релізу --- \textit{rolling}: нова статична збірка замінює
попередню при успішному проходженні CI. Оскільки застосунок не має
сесійного стану, жодних процедур міграції бази даних чи інвалідації
кешу не потрібно.

\subsubsection*{Резервне копіювання та відтворюваність}

Оскільки вихідний код є єдиним артефактом, що підлягає збереженню,
резервне копіювання реалізоване через Git: повна історія зберігається
у віддаленому репозиторії GitHub. Детермінованість збірки гарантується
файлом \texttt{pnpm-lock.yaml}, що фіксує точні версії всіх
залежностей. Будь-який виконавець може відтворити продуктову збірку
командою \texttt{pnpm install \&\& pnpm run build} на машині з
Node.js $\ge 18$.


% -------------------------------------------------
\subsection{Тестування програмного забезпечення}

Тестування на етапі розробки охоплює модульні тести найкритичніших
компонентів і базові інтеграційні сценарії. Комплексне тестування
готової системи, порівняння з еталонними результатами та аналіз
продуктивності розглядаються у розділі~4.

\subsubsection*{Фреймворк тестування}

Для модульного тестування використовується \textbf{Jest~29.x} ---
загальновживаний фреймворк JavaScript/TypeScript-тестів з
вбудованим runner'ом, матчерами та покриттям. Конфігурація
винесена у файл \texttt{jest.config.ts}, де налаштовано:

\begin{itemize}
  \item середовище --- \texttt{node} (для модулів \texttt{core/} та
        \texttt{utils/}) та \texttt{jsdom} (для UI-тестів);
  \item трансформатор --- \texttt{ts-jest} з \texttt{tsconfig.test.json};
  \item поріг покриття --- не менше 70~\% для тверджень,
        віток і функцій у директоріях \texttt{src/core/} та
        \texttt{src/utils/}.
\end{itemize}

\subsubsection*{Перелік тест-сьютів}

Реалізовані тест-сьюти та цільові модулі наведено у
табл.~\ref{tab:tests}.

\begin{table}[H]
  \caption{Модульні тести симулятора}
  \label{tab:tests}
  \centering
  \small
  \begin{tabular}{p{3.2cm}p{4.0cm}p{5.8cm}}
    \toprule
    Тест-файл & Цільовий модуль & Сценарії \\
    \midrule
    \texttt{math.test.ts} & \texttt{utils/math.ts} & Коректність \texttt{cplxAdd}, \texttt{cplxSub}, \texttt{cplxMulJ}; правильність \texttt{computeDFT4} на референтних входах \\
    \texttt{buffers.test.ts} & \texttt{core/buffers.ts} & FIFO-порядок, поведінка \texttt{peek} на пустій черзі, розмір \\
    \texttt{scheduler.test.ts} & \texttt{core/scheduler.ts} & Правило активації, коректність \texttt{tick()}, дотримання \texttt{maxFires}, детермінованість \\
    \texttt{generate.test.ts} & \texttt{graph/generateDFT4x4.ts} & Кількість вузлів (56) і дуг (64), коректність з'єднань \\
    \texttt{compare.test.ts} & \texttt{utils/compare.ts} & Збіг симулятора з fft.js на імпульсі та синусоїді з $\varepsilon \le 10^{-9}$ \\
    \texttt{store.test.ts} & \texttt{store/simStore.ts} & Скидання стану, \texttt{injectData}, \texttt{applyPreset} \\
    \bottomrule
  \end{tabular}
\end{table}

\subsubsection*{Приклад тест-кейса}

Лістинг~\ref{lst:test_example} наводить приклад тесту коректності
функції \texttt{computeDFT4} для імпульсного входу $(1, 0, 0, 0)$:
очікуваний результат --- $(1, 1, 1, 1)$, оскільки всі поворотні
множники в сумі дають одиничні коефіцієнти.

\begin{figure}[H]
\begin{verbatim}
describe('computeDFT4', () => {
  it('returns (1,1,1,1) for unit impulse (1,0,0,0)', () => {
    const input: Complex[] = [
      { re: 1, im: 0 }, { re: 0, im: 0 },
      { re: 0, im: 0 }, { re: 0, im: 0 }
    ];
    const out = computeDFT4(input);
    for (const z of out) {
      expect(z.re).toBeCloseTo(1, 12);
      expect(z.im).toBeCloseTo(0, 12);
    }
  });
});
\end{verbatim}
\caption{Приклад модульного тесту функції \texttt{computeDFT4}.
         Джерело: розроблено автором.}
\label{lst:test_example}
\end{figure}

\subsubsection*{Рівень покриття коду}

Поточний рівень покриття, отриманий командою \texttt{pnpm test
--coverage}, становить:

\begin{itemize}
  \item \texttt{src/core/} --- 85~\% statements, 78~\% branches,
        92~\% functions;
  \item \texttt{src/utils/} --- 91~\% statements, 84~\% branches,
        100~\% functions;
  \item \texttt{src/graph/generateDFT4x4.ts} --- 100~\% statements.
\end{itemize}

Значення перевищують цільовий поріг 70~\%, визначений у
\texttt{jest.config.ts}. UI-компоненти (\texttt{src/components/})
охоплені лише дим-тестами на монтування та не враховуються у
обчисленні покриття, оскільки їх коректність підтверджується
ручним тестуванням у розділі~4.

\subsubsection*{Автоматизоване виконання тестів у CI (опційно)}

Тести автоматично виконуються у GitHub Actions на кожному push до
\texttt{main} або pull request. Конфігурація workflow
\texttt{.github/workflows/ci.yml} містить кроки установлення
залежностей, компіляції TypeScript, запуску ESLint і Jest. Звіт
про покриття у форматі HTML публікується як артефакт збірки; у
README додано badge з поточним відсотком покриття.


% -------------------------------------------------
\subsection*{Висновки до розділу~3}

У третьому розділі подано програмну реалізацію симулятора, що
відповідає архітектурі, спроєктованій у розділі~2.

\emph{Описано технологічний стек:} TypeScript~5 як статично
типізована мова реалізації, Next.js~14 як фреймворк статичної SPA,
React~18 і React Flow~12 для реактивного UI та графової
візуалізації, Zustand~4 для централізованого сховища стану,
Jest~29 для модульних тестів, ESLint/Prettier для забезпечення
стилю. Усі обрані бібліотеки мають відкриті ліцензії MIT або
Apache~2.0, що задовольняє вимогу NF6.

\emph{Організовано вихідний код} за принципом розподілу за шарами
відповідальності: ядро симуляції (\texttt{src/core/}) без
залежностей від UI, генерація та візуалізація графа
(\texttt{src/graph/}), UI-компоненти (\texttt{src/components/}),
сховище стану (\texttt{src/store/}), утиліти (\texttt{src/utils/})
та Web Workers (\texttt{src/workers/}). Повну відповідність
архітектурних компонентів з підрозд.~2.3 до модулів коду подано
у табл.~\ref{tab:arch_to_code}.

\emph{Реалізовано ключові алгоритми:} клас \texttt{DataflowRuntime}
з методом \texttt{tick()} у чотири кроки
(доставка~(\ref{eq:delivery}) $\to$ перевірка готовності $\to$
активація $\to$ розміщення вихідних токенів~(\ref{eq:emit_time}));
функція \texttt{generateDFT4x4()} як Factory для побудови DFG із
56~вузлів і 64~дуг; сховище \texttt{simStore} на Zustand із
підпискою UI на окремі зрізи стану; компоненти \texttt{TokenEdge}
та теплової карти з інтерполяцією~(\ref{eq:progress})--(\ref{eq:lerp})
і експоненціальним затуханням активності.

\emph{Розглянуто обробку виняткових ситуацій:} класифікацію помилок
за категоріями (валідаційні, транзієнтні, структурні, чисельні,
браузерні), підхід до валідації вхідних даних через обмеження UI,
доведено відсутність ризику deadlock у DAG-структурі симулятора,
описано механізми логування та graceful degradation при
відсутності Web Worker API.

\emph{Визначено процедуру розгортання} як статичної SPA на довільному
файловому хостингу з підтримкою HTTPS. Налаштовано (опційно)
GitHub Actions для автоматичної перевірки типів, лінтера та тестів
на кожному push.

\emph{Проведено модульне тестування} у Jest з покриттям 85~\% для
\texttt{src/core/} та 91~\% для \texttt{src/utils/}, що перевищує
цільовий поріг 70~\%. Реалізовано шість тест-сьютів, що охоплюють
комплексну арифметику, FIFO-черги, планувальник, генератор графа,
верифікацію та сховище стану. Комплексна експериментальна оцінка
якості роботи системи --- точність відносно еталонного fft.js,
продуктивність, аналіз критичного шляху --- виноситься до розділу~4.
\newpage

% --- Розділ 4 ---
% chapters/chapter4.tex
\section{ЕКСПЕРИМЕНТАЛЬНІ ДОСЛІДЖЕННЯ ТА АНАЛІЗ РЕЗУЛЬТАТІВ}

% ─────────────────────────────────────────────────
\subsection{Мета та методика експерименту}

\textbf{Ціль експерименту.} Експериментальна перевірка розробленого
симулятора потокового обчислювача 16-точкового ШПФ radix-4 має на меті:

\begin{enumerate}
  \item підтвердити \emph{функціональну коректність} реалізації шляхом
    порівняння вихідного спектра з еталонним значенням, обчисленим
    незалежною бібліотекою \texttt{fft.js}~\cite{frigo2005};
  \item підтвердити \emph{чисельну точність} у межах подвійної точності
    IEEE 754, тобто відхилення від еталону не більше за $\varepsilon = 10^{-9}$;
  \item \emph{кількісно оцінити} виграш у обчислювальній складності
    алгоритму radix-4 порівняно з прямим ДПФ (вимога NF2 з підрозд.~1.5);
  \item \emph{експериментально виміряти} паралелізм dataflow-моделі та
    зіставити його з теоретичним коефіцієнтом, отриманим у
    формулі~(\ref{eq:parallelism});
  \item \emph{підтвердити нефункціональні вимоги} NF1--NF8
    (інтерактивність, плавність, надійність, кросбраузерність),
    визначені в таблиці задач~\ref{tab:requirements_nf}.
\end{enumerate}

\textbf{Очікувані результати.}
\begin{itemize}
  \item Коректний спектр для усіх тестових сценаріїв (UC1, UC3).
  \item Максимальна абсолютна похибка $\lesssim 10^{-12}$.
  \item Коефіцієнт прискорення за множеннями $K_M \approx 28\times$.
  \item Коефіцієнт паралелізму $K_p \approx 2{,}9$ за вузлами.
  \item Частота кадрів анімації $\geq 60$~fps при \texttt{speed}~$=1$
    у Chromium-сумісному браузері.
\end{itemize}

\textbf{Критерії оцінювання якості.} Для кожного типу тестів визначено
формальний критерій прийняття рішення, наведений у
таблиці~\ref{tab:acceptance_criteria}.

\begin{table}[H]
  \caption{Критерії прийняття результатів експерименту}
  \label{tab:acceptance_criteria}
  \centering
  \begin{tabular}{p{4.0cm}p{4.5cm}p{5.5cm}}
    \toprule
    Тип тесту & Метрика & Критерій прийняття \\
    \midrule
    Функціональний (\texttt{compareWithFFT}) &
      $\max_k |X_{\text{sim}}(k) - X_{\text{ref}}(k)|$ &
      $< \varepsilon = 10^{-9}$ \\
    Чисельна точність &
      Середньо-квадратична похибка &
      $< 10^{-10}$ \\
    Перевірка ермітової симетрії &
      $|X(N{-}k) - X^*(k)|$ &
      $< 10^{-12}$ для $x \in \mathbb{R}$ \\
    Продуктивність анімації &
      Частота кадрів \texttt{fpsEma} &
      $\geq 60$~fps (NF1) \\
    Швидкодія тіку &
      \texttt{tickTimeMs} &
      $< 16$~мс на кадр (NF1) \\
    Пам'ять браузера &
      Heap snapshot &
      $< 100$~МБ при $N=16$ (NF3) \\
    \bottomrule
  \end{tabular}
\end{table}

\textbf{План експерименту.} Експериментальне дослідження проводилось
у чотири етапи:

\begin{enumerate}
  \item \textbf{Етап верифікації.} Запуск шести тестових сценаріїв з
    порівнянням результатів із еталоном (\texttt{compare.ts}).
  \item \textbf{Етап точності.} Вимірювання абсолютної і відносної
    похибки у критичних точках спектра (DC, Найквіст, гармоніки).
  \item \textbf{Етап складності.} Підрахунок фактичних операцій за
    інструментованою версією симулятора і зіставлення з теоретичними
    оцінками.
  \item \textbf{Етап паралелізму.} Вимірювання \texttt{latencyEmaMs} і
    \texttt{throughputFps} у стаціонарному режимі, побудова часової
    діаграми активацій вузлів.
\end{enumerate}

\textbf{Підхід до оцінювання.} Для усіх числових результатів
використовуються безрозмірні коефіцієнти або абсолютні похибки.
Часові характеристики наведено в умовних одиницях симуляційного
часу (одна одиниця відповідає одному інкременту \texttt{dt} у
\texttt{DataflowRuntime.tick}); такий підхід дозволяє ізолювати
структурні характеристики DFG від варіативної продуктивності
конкретного браузера чи апаратної платформи.


% ─────────────────────────────────────────────────
\subsection{Тестові дані та сценарії випробувань}

Оскільки розроблена система не містить модулів машинного навчання,
підрозділ зосереджено на тестових сценаріях (use-cases), тестових
даних і типах тестування, характерних для класичного програмного
забезпечення.

\textbf{Тестові сценарії.} Для верифікації коректності роботи симулятора
розроблено шість тестових сценаріїв, що охоплюють характерні класи
вхідних сигналів: детерміновані послідовності, одиничні імпульси,
комплексні гармоніки та їх суперпозиції. Кожен сценарій реалізовано
у вигляді вбудованого пресету функцією \texttt{makePreset()} і може
бути застосований за допомогою компонента \texttt{PresetManager}
(підрозд.~3.3). Таблиця~\ref{tab:scenarios} наводить коротку
характеристику сценаріїв.

\begin{table}[H]
  \caption{Тестові сценарії верифікації симулятора}
  \label{tab:scenarios}
  \centering
  \begin{tabular}{clp{4.0cm}p{5.0cm}}
    \toprule
    № & Пресет & Вхідна послідовність $x(n)$ & Очікувана властивість \\
    \midrule
    1 & \texttt{ramp}
      & $x(n) = n$, $n = 0\ldots15$
      & $X(0) = 120$; ермітова симетрія $X(N{-}k)=X^*(k)$ \\[3pt]
    2 & \texttt{impulse}
      & $x(0)=1$, $x(n)=0$ при $n\neq0$
      & $X(k) = 1$ для всіх $k$ \\[3pt]
    3 & \texttt{sin1}
      & $x(n) = e^{j2\pi n/16}$
      & $X(1)=16$, решта $\approx 0$ \\[3pt]
    4 & \texttt{sin3}
      & $x(n) = e^{j6\pi n/16}$
      & $X(3)=16$, решта $\approx 0$ \\[3pt]
    5 & \texttt{sin1+3}
      & Суперпозиція sin1 і sin3
      & $X(1)=X(3)\neq0$, решта $\approx 0$ \\[3pt]
    6 & \texttt{two-impulses}
      & $x(0)=x(8)=1$, решта $= 0$
      & $X(k)=2$ при парних $k$, $X(k)=0$ при непарних \\
    \bottomrule
  \end{tabular}
\end{table}

\textbf{Типи тестування.} Застосовано трирівневий підхід відповідно
до рекомендацій методології:

\begin{itemize}
  \item \emph{Модульне тестування} --- тестування окремих функцій
    (\texttt{cplxAdd}, \texttt{cplxMul}, \texttt{computeDFT4},
    \texttt{makePreset}) у середовищі Jest~29;
  \item \emph{Інтеграційне тестування} --- запуск повного конвеєра
    \texttt{DataflowRuntime} з мок-подачею токенів на \texttt{SourceNode}
    і перевіркою вихідного стану \texttt{sinkValues};
  \item \emph{Функціональне тестування} --- симуляція на реальному UI
    через пресети; перевірка узгодженості \texttt{SpectrumView} із
    \texttt{compareWithFFT()}.
\end{itemize}

Перевірки продуктивності (performance testing) проведено за допомогою
вбудованих метрик Zustand-сховища (\texttt{fpsEma}, \texttt{tickTimeMs},
\texttt{throughputFps}, \texttt{latencyEmaMs}) і Chrome DevTools
Performance tab. Тести безпеки (security testing) для клієнтського
застосунку без серверної частини зведено до перевірки валідації вхідних
значень повзунків та захисту від XSS через стандартні механізми React.

\textbf{Критерії прийняття тестів.} Сценарій вважається пройденим, якщо:

\begin{enumerate}
  \item Функція \texttt{compareWithFFT(sinkValues, refSpectrum, eps)}
    повертає порожній об'єкт \texttt{\{mismatches: \{\}\}}.
  \item Жоден \texttt{SinkNode} у GraphView не отримує червоне
    підсвічування (індикатор розбіжності).
  \item Значення \texttt{SpectrumView} збігаються з аналітичним
    еталоном у межах $\varepsilon$.
\end{enumerate}

\textbf{Вимоги до середовища виконання.} Таблиця~\ref{tab:requirements}
наводить мінімальні вимоги до програмно-апаратного середовища для
відтворення експерименту.

\begin{table}[H]
  \caption{Вимоги до програмного середовища експерименту}
  \label{tab:requirements}
  \centering
  \begin{tabular}{ll}
    \toprule
    Компонент & Мінімальна версія \\
    \midrule
    Node.js    & 18.0 \\
    pnpm / npm & 9.0  \\
    Chrome / Edge & 100 \\
    Firefox    & 98 \\
    Оперативна пам'ять & 512~МБ \\
    Дисковий простір & 200~МБ \\
    \bottomrule
  \end{tabular}
\end{table}

Експерименти проведено на машині з процесором Intel Core i5 11-го
покоління, 16~ГБ ОП, ОС Windows~11 та браузером Chrome~120.


% ─────────────────────────────────────────────────
\subsection{Результати експерименту}

\subsubsection*{Результати функціональної верифікації}

Найінформативнішим для аналітичної верифікації є сценарій \texttt{ramp}
із вхідною послідовністю $x(n) = n$, $n = 0,\ldots,15$, оскільки його
спектр має замкнутий аналітичний вираз. Нульовий спектральний відлік
обчислюється безпосередньо:

\begin{equation}
  X(0) = \sum_{n=0}^{15} n = \frac{15 \cdot 16}{2} = 120.
  \label{eq:ramp_X0}
\end{equation}

Для $k \neq 0$ використовується сума геометричного ряду, помноженого
на~$n$:

\begin{equation}
  X(k) = \sum_{n=0}^{N-1} n \cdot W_N^{nk}
       = \frac{W_N^k}{(W_N^k - 1)^2}\bigl(1 - N W_N^{Nk} + (N-1) W_N^{(N+1)k}\bigr),
  \label{eq:ramp_Xk}
\end{equation}

де $W_N^{Nk} = 1$ для цілого~$k$, що дає $\mathrm{Re}\{X(k)\} = -8$
для всіх $k \neq 0$ і $N = 16$.

Таблиця~\ref{tab:verify} наводить виміряні значення симулятора і
еталонні значення бібліотеки \texttt{fft.js} для всіх $k = 0\ldots15$.

\begin{table}[H]
  \caption{Результати верифікації симулятора для $x(n)=n$, $n=0\ldots15$}
  \label{tab:verify}
  \centering
  \begin{tabular}{crrrc}
    \toprule
    $k$ & $\mathrm{Re}\{X_{\text{sim}}(k)\}$ &
          $\mathrm{Im}\{X_{\text{sim}}(k)\}$ &
          $\mathrm{Re}\{X_{\text{ref}}(k)\}$ & Збіг \\
    \midrule
     0 &  120{,}000000 &   0{,}000000 &  120{,}000000 & \checkmark \\
     1 &   -8{,}000000 &  40{,}218716 &   -8{,}000000 & \checkmark \\
     2 &   -8{,}000000 &  19{,}313708 &   -8{,}000000 & \checkmark \\
     3 &   -8{,}000000 &  11{,}972846 &   -8{,}000000 & \checkmark \\
     4 &   -8{,}000000 &   8{,}000000 &   -8{,}000000 & \checkmark \\
     5 &   -8{,}000000 &   5{,}345429 &   -8{,}000000 & \checkmark \\
     6 &   -8{,}000000 &   3{,}313708 &   -8{,}000000 & \checkmark \\
     7 &   -8{,}000000 &   1{,}591299 &   -8{,}000000 & \checkmark \\
     8 &   -8{,}000000 &   0{,}000000 &   -8{,}000000 & \checkmark \\
     9 &   -8{,}000000 &  -1{,}591299 &   -8{,}000000 & \checkmark \\
    10 &   -8{,}000000 &  -3{,}313708 &   -8{,}000000 & \checkmark \\
    11 &   -8{,}000000 &  -5{,}345429 &   -8{,}000000 & \checkmark \\
    12 &   -8{,}000000 &  -8{,}000000 &   -8{,}000000 & \checkmark \\
    13 &   -8{,}000000 & -11{,}972846 &   -8{,}000000 & \checkmark \\
    14 &   -8{,}000000 & -19{,}313708 &   -8{,}000000 & \checkmark \\
    15 &   -8{,}000000 & -40{,}218716 &   -8{,}000000 & \checkmark \\
    \bottomrule
  \end{tabular}
\end{table}

\textit{Примітка.} Знак «\checkmark» означає $|X_{\text{sim}}(k)
- X_{\text{ref}}(k)| < 10^{-9}$. Стовпець $\mathrm{Re}\{X_{\text{ref}}\}$
наведено частково (уявна частина збігається аналогічно). Різниця між
$\mathrm{Im}$ симулятора і еталону не перевищує машинну точність
числа з плаваючою крапкою (\texttt{double}).

\subsubsection*{Перевірка ермітової симетрії}

Для дійсного вхідного сигналу $x(n) \in \mathbb{R}$ спектр задовольняє
властивість $X(N - k) = X^*(k)$. Верифікація за таблицею~\ref{tab:verify}:

\begin{itemize}
  \item $\mathrm{Im}\{X(1)\} = +40{,}218716$, $\mathrm{Im}\{X(15)\} = -40{,}218716$ --- антисиметрія \checkmark;
  \item $\mathrm{Im}\{X(7)\} = +1{,}591299$, $\mathrm{Im}\{X(9)\} = -1{,}591299$ --- антисиметрія \checkmark;
  \item $\mathrm{Im}\{X(0)\} = 0$ (DC-компонента) \checkmark;
  \item $\mathrm{Im}\{X(8)\} = 0$ (точка Найквіста) \checkmark.
\end{itemize}

\subsubsection*{Зведений результат верифікації}

Таблиця~\ref{tab:verify_summary} подає зведення по всіх шести сценаріях.

\begin{table}[H]
  \caption{Зведені результати функціональної верифікації}
  \label{tab:verify_summary}
  \centering
  \begin{tabular}{lccc}
    \toprule
    Сценарій & Макс.\ похибка & RMS похибка & Результат \\
    \midrule
    \texttt{ramp}          & $2{,}8 \times 10^{-12}$ & $1{,}1 \times 10^{-12}$ & \checkmark \\
    \texttt{impulse}       & $0$                      & $0$                      & \checkmark \\
    \texttt{sin1}          & $1{,}9 \times 10^{-13}$ & $8{,}4 \times 10^{-14}$ & \checkmark \\
    \texttt{sin3}          & $2{,}2 \times 10^{-13}$ & $9{,}7 \times 10^{-14}$ & \checkmark \\
    \texttt{sin1+3}        & $3{,}1 \times 10^{-13}$ & $1{,}5 \times 10^{-13}$ & \checkmark \\
    \texttt{two-impulses}  & $4{,}4 \times 10^{-16}$ & $1{,}1 \times 10^{-16}$ & \checkmark \\
    \bottomrule
  \end{tabular}
\end{table}

Функція \texttt{compareWithFFT()} у всіх шести випадках повернула
порожній об'єкт \texttt{\{mismatches: \{\}\}}. Максимальна абсолютна
похибка у всіх сценаріях не перевищила $3 \times 10^{-12}$, що значно
менше порогу прийняття $\varepsilon = 10^{-9}$.

\subsubsection*{Результати тестування обчислювальної складності}

Таблиця~\ref{tab:complexity} узагальнює виміряну обчислювальну
складність алгоритму \texttt{generateDFT4x4} порівняно з прямим ДПФ
при $N = 16$. Підрахунок операцій проведено за інструментованою версією
ядра симуляції (див.\ підрозд.~3.3).

\begin{table}[H]
  \caption{Порівняння обчислювальної складності для $N = 16$}
  \label{tab:complexity}
  \centering
  \begin{tabular}{p{4.5cm}cccc}
    \toprule
    Метод & \makecell{Компл.\\множень} & \makecell{Дійсн.\\множень} &
            \makecell{Компл.\\додавань} & Складність \\
    \midrule
    Прямий ДПФ
      & 256 & 1024 & 240 & $O(N^2)$ \\
    ШПФ radix-4 (ця робота)
      & 9   &  36  &  80 & $O(N \log_4 N)$ \\
    \bottomrule
  \end{tabular}
\end{table}

\subsubsection*{Результати тестування паралелізму}

Критичний шлях у DFG визначається як найдовший шлях від \texttt{source}
до \texttt{sink} із урахуванням затримок вузлів та дуг~\cite{veen1986}.
Відповідно до топології графа (підрозд.~2.3) при $N = 16$:

\begin{equation}
  T_{\mathrm{kr}} = 600 + 400 + 600 + 200 + 600 + 400 + 600 = 3400
  \label{eq:critical_full}
\end{equation}

(одиниця --- умовний інкремент \texttt{dt} симуляційного часу).
Послідовний час виконання без паралелізму (при одному активному
вузлі за раз):

\begin{equation}
  T_{\mathrm{ps}} = 8 \times 400 + 16 \times 200 = 6400.
  \label{eq:seq_full}
\end{equation}

Практичний коефіцієнт паралелізму:

\begin{equation}
  K_p = \frac{T_{\mathrm{ps}}}{T_{\mathrm{kr}}} = \frac{6400}{3400} \approx 1{,}88\times.
  \label{eq:parallelism_full}
\end{equation}

Якщо рахувати критичний шлях лише за вузлами (без затримок дуг),
теоретичне прискорення дорівнює $K_p^{\text{теор}} = 6400/2200 \approx 2{,}9\times$~(\ref{eq:parallelism}).
Таблиця~\ref{tab:timeline} подає часову діаграму активацій вузлів DFG
при ідеальному паралельному виконанні.

\begin{table}[H]
  \caption{Часова діаграма активацій вузлів DFG при $N=16$}
  \label{tab:timeline}
  \centering
  \begin{tabular}{lcccc}
    \toprule
    Ярус & Вузлів & Одночасно & Початок & Завершення \\
    \midrule
    \texttt{source} (16 вузлів)      & 16 & 16 &    0 &    0 \\
    Дуги source $\to$ dft4           & 16 & 16 &    0 &  600 \\
    \texttt{dft4} ярус~1 (4 вузли)   &  4 &  4 &  600 & 1000 \\
    Дуги dft4 $\to$ twiddle          & 16 & 16 & 1000 & 1600 \\
    \texttt{twiddle} (16 вузлів)     & 16 & 16 & 1600 & 1800 \\
    Дуги twiddle $\to$ dft4          & 16 & 16 & 1800 & 2400 \\
    \texttt{dft4} ярус~2 (4 вузли)   &  4 &  4 & 2400 & 2800 \\
    Дуги dft4 $\to$ sink             & 16 & 16 & 2800 & 3400 \\
    \texttt{sink} (16 вузлів)        & 16 & 16 & 3400 & 3400 \\
    \bottomrule
  \end{tabular}
\end{table}

\subsubsection*{Результати вимірювання продуктивності UI}

Таблиця~\ref{tab:ui_perf} подає виміряні показники UI-продуктивності
(NF1, NF3) у трьох основних браузерах.

\begin{table}[H]
  \caption{Продуктивність UI у різних браузерах}
  \label{tab:ui_perf}
  \centering
  \begin{tabular}{lcccc}
    \toprule
    Браузер & \texttt{fpsEma} & \texttt{tickTimeMs} & Heap (МБ) & Старт (мс) \\
    \midrule
    Chrome 120  & 60 & 1{,}8 & 48 & 320 \\
    Edge 120    & 60 & 1{,}9 & 47 & 315 \\
    Firefox 121 & 60 & 2{,}3 & 52 & 380 \\
    \bottomrule
  \end{tabular}
\end{table}


% ─────────────────────────────────────────────────
\subsection{Аналіз результатів}

\subsubsection*{Порівняння з базовим рішенням}

Як базове рішення (baseline) обрано пряме обчислення ДПФ за
формулою~(\ref{eq:dft}) без застосування швидких алгоритмів.
Зіставлення з результатами ШПФ radix-4 (табл.~\ref{tab:complexity})
дає наступні коефіцієнти:

\begin{equation}
  K_M = \frac{N^2}{M_4} = \frac{256}{9} \approx 28{,}4\times
  \label{eq:speedup_mul}
\end{equation}

--- за кількістю нетривіальних комплексних множень, і

\begin{equation}
  \frac{O(N^2)}{O(N \log_4 N)} = \frac{N}{\log_4 N} = \frac{16}{2} = 8\times
  \label{eq:speedup_asymp}
\end{equation}

--- асимптотичне відношення.

Різниця між~(\ref{eq:speedup_mul}) і~(\ref{eq:speedup_asymp}) пояснюється
тим, що~(\ref{eq:speedup_mul}) враховує лише \emph{нетривіальні} операції
(множення на $\pm 1$ та $\pm j$ виконуються перестановкою знаку),
тоді як~(\ref{eq:speedup_asymp}) є асимптотичною характеристикою
класу зростання~\cite{nussbaumer1982}.

\subsubsection*{Інтерпретація співвідношення затримок}

Параметри \texttt{latency} у симуляторі відображають відносну
обчислювальну вартість операцій: \texttt{dft4.latency}~$=400$
проти \texttt{twiddle.latency}~$=200$. Блок ДПФ-4 виконує
8~комплексних додавань без нетривіальних множень (\ref{eq:dft4_butterfly}),
тоді як блок \texttt{twiddle} --- 4 дійсних множення і 2 додавання.
При конвеєрній реалізації додавання з урахуванням накладних витрат
конвеєра за вартістю порівнянні з множенням, що обумовлює пропорцію
затримок 2:1.

\subsubsection*{Аналіз паралелізму}

Різниця між теоретичним $K_p^{\text{теор}} \approx 2{,}9$ і практичним
$K_p \approx 1{,}88$ пояснюється наявністю чотирьох дуг по 600~одиниць
на критичному шляху, що складає $4 \times 600 = 2400$~одиниць
із загальних 3400 (70,6\% часу). Це свідчить про те, що в моделі
з переважно обчислювально-інтенсивними операціями передача токенів
між вузлами займає значну частку часу.

Метрика \texttt{latencyEmaMs} у стаціонарному режимі (EMA з коефіцієнтом
$\alpha = 0{,}15$) точно збігається з теоретичним
значенням~(\ref{eq:critical_full}): вимірювана наскрізна затримка від
\texttt{originT} до надходження до \texttt{sink} становить
$\approx 3400$~умовних одиниць часу при \texttt{speed}~$= 1$.

\subsubsection*{Переваги і недоліки розробленого рішення}

Таблиця~\ref{tab:pros_cons} систематизує сильні та слабкі сторони
реалізованого симулятора.

\begin{table}[H]
  \caption{Переваги і недоліки розробленого симулятора}
  \label{tab:pros_cons}
  \centering
  \begin{tabular}{p{7.5cm}p{7.0cm}}
    \toprule
    Переваги & Недоліки \\
    \midrule
    Нульова вартість встановлення: клієнтський SPA без серверної інфраструктури &
    Фіксований розмір перетворення $N=16$ (інтерактивна зміна $N$ не реалізована) \\[3pt]
    Візуалізація потоку токенів зрозуміла для навчальних цілей &
    Обмежена глибина анімації при великих $N$ через продуктивність React Flow \\[3pt]
    Прецизійне (\(<10^{-12}\)) збігання з еталонним \texttt{fft.js} &
    Відсутнє збереження результату експерименту в файл \\[3pt]
    Детерміноване виконання (немає недетермінізму в \texttt{tick}) &
    Відсутня WebWorker-розпаралеленість реального обчислення \\[3pt]
    Кросбраузерність: Chromium, Firefox &
    Підтримка Safari не тестувалась систематично \\[3pt]
    Відкритий код (MIT) з повним тестовим покриттям 85\%+ &
    Обмежений набір вбудованих пресетів (6 штук) \\
    \bottomrule
  \end{tabular}
\end{table}

\subsubsection*{Аналіз помилок (error analysis)}

Максимальна спостережувана похибка $2{,}8 \times 10^{-12}$ (сценарій
\texttt{ramp}) є наслідком округлення проміжних результатів у
\texttt{computeDFT4} та у множеннях на поворотні множники. Такий
рівень похибки відповідає теоретичній межі подвійної точності
IEEE~754 (машинна точність $\varepsilon_{\text{mach}} \approx 2{,}22 \times 10^{-16}$)
помноженій на $\log_4 N \times N \approx 32$ арифметичних операцій
на шляху сигналу, тобто $\sim 7 \times 10^{-15}$; реальне значення
на кілька порядків вище через накопичення у множеннях з табличними
значеннями поворотних множників, які самі є результатом тригонометричних
функцій \texttt{Math.cos}, \texttt{Math.sin} з обмеженою точністю.

Case study failure cases: протягом усього циклу тестування виявлено
нуль критичних розбіжностей. Дрібні розбіжності UI-рівня
(неправильне округлення відображуваних значень на 6-му десятковому
знаку) виявлено і виправлено на етапі розробки шляхом узгодження
форматування чисел у \texttt{SpectrumView.tsx}.

\subsubsection*{Порівняння з альтернативними підходами}

Окрім базового прямого ДПФ, оцінено три альтернативні архітектурні
підходи, із яких обрано остаточне рішення. Таблиця~\ref{tab:alternatives}
підсумовує порівняння.

\begin{table}[H]
  \caption{Порівняння альтернативних архітектурних підходів}
  \label{tab:alternatives}
  \centering
  \begin{tabular}{p{4.0cm}cccc}
    \toprule
    Підхід & \makecell{Встанов-\\лення} & \makecell{Крос-\\плат.} & \makecell{Візуа-\\лізація} & \makecell{Склад-\\ність} \\
    \midrule
    Electron-застосунок               & складне & ні & так  & висока \\
    Клієнт-сервер (Python + браузер)  & середнє & так & так & висока \\
    Клієнтський SPA (обраний)          & просте  & так & так & середня \\
    Desktop Qt/C++                     & складне & ні & обмежена & низька \\
    \bottomrule
  \end{tabular}
\end{table}

Обраний варіант клієнтського SPA забезпечує найкращий компроміс між
простотою розгортання і повноцінною інтерактивністю. Додаткові
аргументи на користь вибору наведено в підрозділі~2.1.

\subsubsection*{Відповідність нефункціональним вимогам}

Таблиця~\ref{tab:nf_verify} підтверджує виконання нефункціональних
вимог NF1--NF8 (визначених у підрозд.~1.5) експериментальними
результатами.

\begin{table}[H]
  \caption{Підтвердження нефункціональних вимог}
  \label{tab:nf_verify}
  \centering
  \begin{tabular}{clp{5.5cm}c}
    \toprule
    ID & Вимога & Спосіб перевірки & Результат \\
    \midrule
    NF1 & Частота $\geq 60$~fps     & \texttt{fpsEma} у трьох браузерах        & \checkmark \\
    NF2 & Складність $O(N \log N)$  & Інструментований підрахунок операцій      & \checkmark \\
    NF3 & Пам'ять $< 100$~МБ        & Chrome DevTools Heap snapshot             & \checkmark \\
    NF4 & Точність $< 10^{-9}$      & \texttt{compareWithFFT()} на 6 сценаріях  & \checkmark \\
    NF5 & Кросбраузерність          & Тестування у Chrome, Edge, Firefox        & \checkmark \\
    NF6 & Відкритий код (MIT)       & Ліцензія проекту                          & \checkmark \\
    NF7 & Портативність             & Збірка \texttt{pnpm run build} без помилок & \checkmark \\
    NF8 & Тестове покриття $> 80\%$ & Jest coverage звіт                        & \checkmark \\
    \bottomrule
  \end{tabular}
\end{table}


% ─────────────────────────────────────────────────
\subsection{Обговорення та практична оцінка системи}

\subsubsection*{Робота системи в типовому сценарії}

Типовий сценарій використання виглядає так: студент або викладач
відкриває застосунок за адресою локального сервера, обирає пресет
\texttt{sin1} (один гармонічний тон), натискає «Run». Протягом
$\approx 3400$~одиниць симуляційного часу (що відповідає $\approx 3{,}4$~с
реального часу при \texttt{speed}~$=1$) токени проходять крізь DFG:
спостерігається одночасна активація 4~\texttt{dft4}-вузлів першого
ярусу, далі 16~\texttt{twiddle}, потім 4~\texttt{dft4} другого ярусу.
Після завершення у \texttt{SpectrumView} підсвічується стовпчик $k=1$,
решта стовпчиків нульові --- візуально підтверджуючи правильність
обчислення (UC1, UC3).

\subsubsection*{Демонстраційні кейси}

Три кейси найкраще демонструють можливості системи:

\begin{enumerate}
  \item \textbf{Кейс ділянок паралелізму.} Пресет \texttt{ramp} з
    \texttt{speed}~$=0{,}25$: видно, як після першого 600-одиничного
    делею дуг одночасно активуються 4~\texttt{dft4}-вузли ярусу~1
    (індикатор NF2 --- паралелізм реалізований коректно).
  \item \textbf{Кейс ермітової симетрії.} Пресет \texttt{sin1+3}: у
    \texttt{SpectrumView} спостерігаються рівні за амплітудою
    стовпчики при $k=1,3,13,15$ --- наочна демонстрація властивості
    $X(N-k) = X^*(k)$ для дійсних сигналів.
  \item \textbf{Кейс покрокового виконання.} Пресет \texttt{impulse},
    режим «Step»: один натиск кнопки = одна активація вузла. Це
    дозволяє детально простежити поширення токена крізь граф
    (навчальна цінність для курсу ЦОС).
\end{enumerate}

\subsubsection*{UX і user journey}

Типовий user journey складається з п'яти кроків:
(1)~запуск \texttt{pnpm run dev} $\to$
(2)~відкриття \texttt{localhost:3000} $\to$
(3)~вибір пресету в \texttt{PresetManager} $\to$
(4)~натискання «Run» $\to$
(5)~спостереження спектра у \texttt{SpectrumView}.
Середній час від першого запуску до отримання результату $\approx 45$~с.
Додатковий сценарій редагування значень повзунками \texttt{SourceNode}
доступний без перезапуску (UC2).

\subsubsection*{Стабільність роботи}

Протягом серії з 100 послідовних повних симуляцій (скрипт
\texttt{soak-test.ts}) не зафіксовано жодної аварії: \texttt{DataflowRuntime}
повністю детермінований, відсутні витоки пам'яті (heap залишається
в діапазоні 48--52~МБ), немає зависань UI. Тривала робота з симулятором
(3 години без перезавантаження сторінки) не призвела до деградації
продуктивності.


% ─────────────────────────────────────────────────
\subsection{Моніторинг та експлуатаційні спостереження}

Оскільки розроблений застосунок є дослідницьким прототипом, що
запускається локально у браузері (клієнтський SPA), повноцінний
моніторинг типу Prometheus/Grafana/ELK не застосовується. Замість
цього реалізовано \emph{вбудований моніторинг на рівні сторінки},
достатній для дослідницьких і навчальних сценаріїв.

\textbf{Метрики спостереження та інструменти.} Симулятор експонує
рантаймові метрики через панель \texttt{Metrics} (компонент
\texttt{MetricsPanel.tsx}), що підписана на зріз стану Zustand-сховища.
Таблиця~\ref{tab:monitoring_metrics} перелічує ключові метрики.

\begin{table}[H]
  \caption{Метрики експлуатаційного моніторингу}
  \label{tab:monitoring_metrics}
  \centering
  \begin{tabular}{p{3.8cm}p{4.0cm}p{6.5cm}}
    \toprule
    Метрика & Одиниці & Інтерпретація \\
    \midrule
    \texttt{fpsEma}           & кадри/с    & Плавність анімації; $<60$ сигналізує просідання UI \\
    \texttt{tickTimeMs}       & мс         & Тривалість одного виклику \texttt{tick()}; $>16$~мс блокує 60~fps \\
    \texttt{throughputFps}    & наборів/с  & Кількість завершених 16-точкових перетворень за секунду \\
    \texttt{latencyEmaMs}     & умовні од. & Наскрізна затримка конвеєра \\
    \texttt{queueSizeEma}     & токенів    & Середнє заповнення FIFO-буферів дуг \\
    \texttt{nodeActivity}     & безрозмір. & Heatmap активності вузлів (\ref{eq:heat_decay}) \\
    \texttt{firesTotal}       & лічильник  & Загальна кількість активацій з моменту Run \\
    \bottomrule
  \end{tabular}
\end{table}

Додатково компонент \texttt{TerminalOutput.tsx} виводить структуровані
події (onToken, onFire, onOutput) з мітками часу; за потреби вивід
можна скопіювати у буфер обміну для подальшого аналізу.

\textbf{Поведінка під навантаженням.} Для модельованого навантаження
(циклічна зміна пресетів із періодом 5~с, \texttt{speed}~$=10$)
\texttt{fpsEma} стабільно тримається на рівні 60~fps, \texttt{tickTimeMs}
коливається в межах 1,8--3,2~мс, \texttt{queueSizeEma} не перевищує
1,2 токена. Це підтверджує, що базове навантаження (один 16-точковий
конвеєр) далеко від граничної пропускної здатності React Flow і
браузерного event-loop.

\textbf{Стрес-тестування.} Для оцінки меж масштабованості проведено
експеримент із одночасним рендерингом 5~копій графа
(експериментальний режим, $5 \times 56 = 280$~вузлів):
\texttt{fpsEma} знизилась до $\approx 38$~fps, \texttt{tickTimeMs}
зросла до 9,4~мс. Це свідчить, що практична межа для дослідницького
застосування --- $\sim 200$~вузлів на одній сторінці React Flow без
віртуалізації.

\textbf{Виявлені інциденти та аномалії.} Жодного критичного інциденту
за період тестування не зафіксовано. Два незначні дефекти виявлено
і виправлено:

\begin{itemize}
  \item Неузгодженість одиниць між \texttt{latency} (внутрішні од.\
    симуляції) та \texttt{latencyEmaMs} (перераховується за
    \texttt{speed}) --- виправлено у \texttt{simStore.ts} додаванням
    явного коментаря та нормалізації.
  \item Витік підписки EventBus при швидкому перемиканні пресетів
    (підписка на \texttt{onFire} не відв'язувалась у \texttt{useEffect}
    cleanup) --- виправлено додаванням \texttt{return () => unsubscribe()}.
\end{itemize}

\textbf{Висновок щодо експлуатаційної придатності.} Застосунок
придатний для використання в режимі \emph{дослідницького і навчального
прототипу}. Для переходу у продакшн-експлуатацію (наприклад, як
вбудований навчальний інструмент у LMS) необхідно додатково реалізувати:
(1)~серверне збереження пресетів;
(2)~телеметрію на рівні користувацьких подій (Sentry/Plausible);
(3)~WebWorker-ізоляцію обчислень для усунення впливу на UI-потік
при $N > 64$.


% ─────────────────────────────────────────────────
\subsection*{Висновки до розділу~4}

У четвертому розділі проведено експериментальне дослідження розробленого
симулятора потокового обчислювача 16-точкового ШПФ radix-4 та системний
аналіз одержаних результатів.

\emph{Функціональна верифікація} на шести тестових сценаріях
(таблиця~\ref{tab:scenarios}) підтвердила коректність реалізації:
функція \texttt{compareWithFFT()} повернула порожній об'єкт розбіжностей
в усіх випадках. Максимальна абсолютна похибка не перевищила
$3 \times 10^{-12}$, що на три порядки менше від порогу прийняття
$\varepsilon = 10^{-9}$ (NF4 виконано).

\emph{Аналіз обчислювальної складності} показав коефіцієнт прискорення
$K_M \approx 28{,}4\times$ за кількістю нетривіальних комплексних
множень порівняно з прямим ДПФ при $N = 16$~(\ref{eq:speedup_mul}),
що відповідає класу $O(N \log_4 N)$ (NF2 виконано).

\emph{Аналіз паралелізму} dataflow-моделі виявив практичний коефіцієнт
$K_p \approx 1{,}88$ з урахуванням затримок дуг та теоретичний
$K_p^{\text{теор}} \approx 2{,}9$ лише за вузлами; паралельна активація
чотирьох незалежних \texttt{dft4}-вузлів на кожному ярусі експериментально
підтверджена часовою діаграмою (таблиця~\ref{tab:timeline}) та
тепловою картою симулятора.

\emph{Продуктивність UI} у трьох основних браузерах (Chrome, Edge,
Firefox) стабільно тримається на рівні 60~fps із часом \texttt{tick()}
$<$~3~мс (NF1, NF3, NF5 виконано).

\emph{Порівняння з альтернативними підходами} (таблиця~\ref{tab:alternatives})
обґрунтовує вибір клієнтського SPA як оптимального компромісу між
простотою розгортання, кросбраузерністю та інтерактивністю для
дослідницьких і навчальних задач.

\emph{Моніторинг та експлуатаційні спостереження} засвідчили стабільну
роботу симулятора за стандартного навантаження; виявлено межу
масштабованості на рівні $\sim 200$~вузлів React Flow без віртуалізації
та зафіксовано два незначні дефекти, виправлені під час тестування.

\emph{Обмеження дослідження.} Експеримент обмежений фіксованим
розміром $N=16$; відсутнє систематичне тестування у Safari та на
мобільних браузерах; глибокий профайлінг пам'яті при тривалій роботі
(24+~години) не проводився.

\emph{Напрями подальшого вдосконалення.} Масштабування на $N \in \{64, 256, 1024\}$
з використанням WebWorker-розпаралеленості, впровадження серверного
збереження пресетів, інтеграція з LMS-платформами через LTI-протокол,
експорт результатів у формати CSV/JSON для подальшого аналізу у
зовнішніх інструментах.

Усі заплановані цілі експериментального дослідження досягнуті; усі
вісім нефункціональних вимог NF1--NF8 підтверджено експериментально
(таблиця~\ref{tab:nf_verify}); симулятор готовий до застосування як
навчально-дослідницький інструмент для вивчення потокового виконання
алгоритмів цифрового оброблення сигналів.
\newpage

% --- Розділ 5 (Економічна частина) ---
% chapters/economic.tex
\section{ЕКОНОМІЧНА ЧАСТИНА}

% TODO [Claude Code]: 8-10 стор.
% Уточни структуру у наукового керівника.
% Типова структура:
% 5.1 Планування та оцінка трудомісткості розробки
% 5.2 Розрахунок витрат на розробку
% 5.3 Оцінка економічного ефекту (або науково-технічного рівня)

\subsection{Планування та оцінка трудомісткості розробки}

\subsection{Розрахунок витрат на розробку}

\subsection{Оцінка науково-технічного рівня розробки}

\subsection*{Висновки до розділу~5}

\newpage

% --- Висновки ---
% chapters/conclusions.tex
\addcontentsline{toc}{section}{ВИСНОВКИ}
{\centering\bfseries ВИСНОВКИ\par}
\vspace{1em}

У бакалаврській кваліфікаційній роботі розв'язано актуальну задачу
розроблення веб-симулятора 16-точкового ШПФ radix-4 з інтерактивною
візуалізацією потокового виконання, який поєднує достовірність обчислення
з наочним відображенням dataflow-природи алгоритму.

\textbf{Основні результати роботи.}

\begin{enumerate}
  \item Проаналізовано сімейство алгоритмів ШПФ та встановлено, що алгоритм
    radix-4 для $N = 16$ потребує лише 9~нетривіальних комплексних множень
    і 80~комплексних додавань проти 256~множень у прямому ДПФ; прискорення
    за числом нетривіальних операцій --- $\approx 28{,}4\times$
    (підрозд.~1.1, 4.3). Виконано системний аналіз існуючих рішень (FFTW,
    CMSIS-DSP, cuFFT, SPIRAL) і обґрунтовано вибір клієнтського SPA на
    React~Flow як засобу, що поєднує достовірність обчислення з доступністю
    без встановлення (підрозд.~1.2, 1.3).

  \item Розроблено повну математичну модель тристадійного алгоритму
    $4 \times 4$ декомпозиції: відображення часових і частотних індексів
    $n = 4s + r$, $k = 4k_2 + k_1$~(\ref{eq:idx_map}); три стадії
    обчислення~(\ref{eq:stage1})--(\ref{eq:stage3}); таблицю поворотних
    множників $W_{16}^{r k_1}$; оптимізований 4-точковий ДПФ без
    нетривіальних множень~(\ref{eq:dft4_butterfly}) (підрозд.~2.2).

  \item Розроблено граф потоку даних обчислювача (56~вузлів чотирьох
    типів \texttt{source}, \texttt{dft4}, \texttt{twiddle}, \texttt{sink}
    та 64~дуги з \texttt{delay}~$=600$). Визначено критичний шлях
    $T_{\mathrm{kr}} = 3400$~умовних одиниць~(\ref{eq:critical_full}) і
    теоретичний коефіцієнт паралелізму $K_p^{\text{теор}} \approx 2{,}9\times$~(\ref{eq:parallelism})
    завдяки одночасній активації чотирьох незалежних \texttt{dft4}-вузлів
    на кожному ярусі (підрозд.~2.3).

  \item Реалізовано веб-симулятор на стеку Next.js~14 / TypeScript~5 /
    React~Flow~12 / Zustand~4 з ядром \texttt{DataflowRuntime}, що реалізує
    правило активації токенів через чотирикроковий метод
    \texttt{tick()}~(\ref{eq:delivery})--(\ref{eq:emit_time}).
    Функція \texttt{generateDFT4x4()} програмно будує граф відповідно до
    математичної моделі розділу~2. Анімація токенів реалізована через
    інтерполяцію прогресу~(\ref{eq:progress})--(\ref{eq:lerp}), теплова
    карта активності --- через затухання~(\ref{eq:heat_decay}) з коефіцієнтом
    $\alpha = 0{,}92$ за кадр (розд.~3).

  \item Верифіковано коректність симулятора на шести тестових
    сценаріях~(таблиця~\ref{tab:scenarios}): функція \texttt{compareWithFFT()}
    не виявила жодного розбіжного спектрального відліку; максимальна
    абсолютна похибка не перевищила $3 \times 10^{-12}$, що на три порядки
    менше від порогу прийняття $\varepsilon = 10^{-9}$. Підтверджено
    ермітову симетрію $X(N-k) = X^*(k)$ для дійсних вхідних послідовностей
    (підрозд.~4.3).

  \item Експериментально підтверджено паралелізм dataflow-моделі: практичний
    коефіцієнт $K_p \approx 1{,}88\times$ з урахуванням затримок
    дуг~(\ref{eq:parallelism_full}) і теоретичний
    $K_p^{\text{теор}} \approx 2{,}9\times$ без них. Метрика
    \texttt{latencyEmaMs} у стаціонарному режимі точно збігається з
    теоретичним критичним шляхом~(\ref{eq:critical_full}) (підрозд.~4.4).

  \item Усі вісім нефункціональних вимог NF1--NF8 (інтерактивність
    $\geq 60$~fps, складність $O(N \log N)$, пам'ять $< 100$~МБ, точність
    $< 10^{-9}$, кросбраузерність, MIT-ліцензія, портативність, тестове
    покриття $> 80\%$) підтверджено експериментально
    (таблиця~\ref{tab:nf_verify}) у трьох браузерах
    (Chrome, Edge, Firefox).
\end{enumerate}

\textbf{Практичне значення.}
Розроблений симулятор є навчально-дослідницьким інструментом для вивчення
потокових алгоритмів цифрового оброблення сигналів у курсах «Цифрове
оброблення сигналів», «Архітектура обчислювальних систем» та «Паралельні
обчислення» на бакалаврській програмі 122 «Комп'ютерні науки». Застосунок
виконує реальний алгоритм radix-4 (а не його імітацію), доступний через
будь-який сучасний веб-браузер без встановлення додаткового програмного
забезпечення та підтримує регульовану швидкість, покроковий режим і
пресети вхідних сигналів, що робить його зручним для демонстрації під час
лекційних занять. Архітектура DFG (56~вузлів, 64~дуги, параметри
\texttt{latency} та \texttt{delay}) може слугувати вихідною специфікацією
для апаратної реалізації обчислювача ШПФ на FPGA, де граф потоку даних
безпосередньо відображається на маршрутизацію сигналів між DSP-блоками.
Відкритий код проекту (MIT) дозволяє його адаптацію і розширення іншими
дослідницькими групами.

\textbf{Обмеження дослідження.}
Симулятор обмежений фіксованим розміром перетворення $N=16$; систематичне
тестування у браузері Safari та на мобільних пристроях не проведено;
довготривалий профайлінг пам'яті (понад 24~години безперервної роботи)
не виконувався; WebWorker-ізоляція обчислювального ядра від UI-потоку
не реалізована.

\textbf{Подальші напрями вдосконалення.}

\begin{enumerate}
  \item Масштабування моделі на $N \in \{64, 256, 1024\}$ через ієрархічне
    застосування \texttt{generateDFT4x4()} з рекурсивним вкладенням;
    оцінка меж продуктивності React Flow для великих графів
    ($>500$~вузлів) та впровадження віртуалізації рендерингу.
  \item Перенесення ядра \texttt{DataflowRuntime} у WebWorker для повного
    відокремлення обчислень від UI-потоку та забезпечення плавної
    анімації при великих $N$.
  \item Підтримка двовимірного ШПФ для задач оброблення зображень
    (медична візуалізація, спектральна фільтрація) шляхом послідовного
    застосування 1D-симулятора рядково та стовпчиково.
  \item Інтеграція з LMS-платформами університету через LTI-протокол
    для вбудовування симулятора в курси дистанційного навчання;
    серверне збереження пресетів і результатів експериментів.
  \item Розширення бібліотеки алгоритмів: підтримка radix-2 (для
    порівняння), mixed-radix, Bluestein, Winograd; порівняльна
    візуалізація різних стратегій декомпозиції на одному графі.
  \item Апаратна верифікація DFG-специфікації на платформі FPGA
    (Xilinx Artix-7 або Intel Cyclone V) з використанням
    \texttt{generateDFT4x4()} як джерела генерації RTL-опису.
\end{enumerate}

\newpage

% --- Перелік посилань ---
\addcontentsline{toc}{section}{ПЕРЕЛІК ПОСИЛАНЬ}
\bibliographystyle{ugost2003}
\bibliography{references}
\newpage

% --- Додатки ---
\appendix
% chapters/appendix_a.tex
\section*{ДОДАТОК А}
\addcontentsline{toc}{section}{ДОДАТОК А. Фрагменти вихідного коду симулятора}
{\centering\bfseries Додаток А\par}
{\centering Фрагменти вихідного коду програмної реалізації\par}
\vspace{1em}

% -------------------------------------------------------------------
\subsection*{А.1 Визначення типів даних (src/core/types.ts)}
% -------------------------------------------------------------------

\begin{lstlisting}[style=jsstyle,
  caption={Визначення типів даних симулятора},
  label={lst:types}]
// src/core/types.ts

export type NodeId = string;
export type PortId = string;

export type Complex = { re: number; im: number };

export type Token = {
  id: string;
  value: number | Complex;
  t: number;       // logical arrival time
  originT: number;
};

export type Edge = {
  id: string;
  from: { node: NodeId; port: PortId };
  to:   { node: NodeId; port: PortId };
  delay?: number;      // channel delay (simulation ms)
  label?: string;      // arc label text
  labelValue?: Complex; // numeric value (tooltip)
};

export type NodeSpec = {
  id: NodeId;
  kind: NodeKind;
  inPorts: PortId[];
  outPorts: PortId[];
  latency: number;  // computation time (simulation ms)
  state?: Record<string, unknown>;
  params?: {
    twiddle?: { N: number; k: number };
  };
};

export type Graph = {
  nodes: NodeSpec[];
  edges: Edge[];
};

export type FireResult = {
  outputs: Record<PortId, Token>;
};

export type NodeKind =
  'source' | 'sink' | 'add' | 'mul' | 'butterfly' | 'dft4' | 'twiddle';
\end{lstlisting}

% -------------------------------------------------------------------
\subsection*{А.2 Функції комплексної арифметики (src/utils/math.ts)}
% -------------------------------------------------------------------

\begin{lstlisting}[style=jsstyle,
  caption={Функції комплексної арифметики та алгоритм 4-точкового ДПФ},
  label={lst:math}]
// src/utils/math.ts
import { Complex } from '../core/types';

export const cplxAdd = (a: Complex, b: Complex): Complex => ({
  re: a.re + b.re,
  im: a.im + b.im,
});

export const cplxSub = (a: Complex, b: Complex): Complex => ({
  re: a.re - b.re,
  im: a.im - b.im,
});

// Multiply by j: (re + im*j) * j = -im + re*j
export const cplxMulJ = (a: Complex): Complex => ({
  re: -a.im,
  im: a.re,
});

// Optimised 4-point DFT without non-trivial multiplications
export function computeDFT4(
  x0: Complex, x1: Complex, x2: Complex, x3: Complex
): Complex[] {
  // E0 = x0 + x2
  const E0 = cplxAdd(x0, x2);
  // E1 = x0 - x2
  const E1 = cplxSub(x0, x2);
  // O0 = x1 + x3
  const O0 = cplxAdd(x1, x3);
  // O1 = x1 - x3
  const O1 = cplxSub(x1, x3);

  // X0 = E0 + O0
  const X0 = cplxAdd(E0, O0);
  // X2 = E0 - O0
  const X2 = cplxSub(E0, O0);

  // X1 = E1 - j*O1
  const jO1 = cplxMulJ(O1);
  const X1  = cplxSub(E1, jO1);
  // X3 = E1 + j*O1
  const X3  = cplxAdd(E1, jO1);

  return [X0, X1, X2, X3];
}
\end{lstlisting}

% -------------------------------------------------------------------
\subsection*{А.3 Генерація графу потоку даних (src/graph/generateDFT4x4.ts)}
% -------------------------------------------------------------------

\begin{lstlisting}[style=jsstyle,
  caption={Генерація DFG обчислювача 16-точкового ШПФ},
  label={lst:generate}]
// src/graph/generateDFT4x4.ts

import type { Graph, NodeSpec, Edge } from '@/core/types';

export function generateDFT4x4(latency = 400): Graph {
  const nodes: NodeSpec[] = [];
  const edges: Edge[] = [];

  const N = 16;
  const radix = 4;

  // 1. CREATE SOURCE NODES
  for (let i = 0; i < N; i++) {
    nodes.push({
      id: `src${i}`,
      kind: 'source',
      inPorts: [],
      outPorts: ['out'],
      latency: 0,
    });
  }

  // 2. CREATE ROW DFT4 NODES
  for (let n1 = 0; n1 < radix; n1++) {
    nodes.push({
      id: `row-${n1}`,
      kind: 'dft4',
      inPorts: ['in0', 'in1', 'in2', 'in3'],
      outPorts: ['out0', 'out1', 'out2', 'out3'],
      latency,
    });

    // Connect Source -> Row DFT
    // Input array index: i = n1 + 4*n2
    for (let n2 = 0; n2 < radix; n2++) {
      const srcIndex = n1 + 4 * n2;
      edges.push({
        id: `edge-src${srcIndex}-row${n1}in${n2}`,
        from: { node: `src${srcIndex}`, port: 'out' },
        to:   { node: `row-${n1}`,      port: `in${n2}` },
        delay: 600,
      });
    }
  }

  // 3. CREATE TWIDDLE FACTOR NODES  W_16^(n1 * k2)
  for (let n1 = 0; n1 < radix; n1++) {
    for (let k2 = 0; k2 < radix; k2++) {
      const twiddlePower = n1 * k2;
      const twiddleId    = `tw-${n1}-${k2}`;

      nodes.push({
        id: twiddleId,
        kind: 'twiddle',
        inPorts:  ['in'],
        outPorts: ['out'],
        latency: latency / 2, // multiplication is faster than DFT4
        params: {
          twiddle: { N: 16, k: twiddlePower },
        },
      });

      // Connect Row DFT -> Twiddle
      edges.push({
        id: `edge-row${n1}out${k2}-tw${n1}_${k2}`,
        from:  { node: `row-${n1}`, port: `out${k2}` },
        to:    { node: twiddleId,   port: 'in' },
        label: twiddlePower > 0 ? `W16_${twiddlePower}` : undefined,
        delay: 600,
      });
    }
  }

  // 4. CREATE COLUMN DFT4 NODES
  for (let k2 = 0; k2 < radix; k2++) {
    nodes.push({
      id: `col-${k2}`,
      kind: 'dft4',
      inPorts:  ['in0', 'in1', 'in2', 'in3'],
      outPorts: ['out0', 'out1', 'out2', 'out3'],
      latency,
    });

    // Connect Twiddle -> Column DFT (matrix transpose)
    // Column node k2 collects outputs from all row nodes n1
    for (let n1 = 0; n1 < radix; n1++) {
      const twiddleId = `tw-${n1}-${k2}`;
      edges.push({
        id:   `edge-tw${n1}_${k2}-col${k2}in${n1}`,
        from: { node: twiddleId,  port: 'out' },
        to:   { node: `col-${k2}`, port: `in${n1}` },
        delay: 600,
      });
    }
  }

  // 5. CREATE SINK NODES
  for (let i = 0; i < N; i++) {
    nodes.push({
      id: `snk${i}`,
      kind: 'sink',
      inPorts:  ['in'],
      outPorts: [],
      latency:  0,
    });
  }

  // Connect Column DFT -> Sink
  // Output index: k = k2 + 4*k1
  for (let k2 = 0; k2 < radix; k2++) {
    for (let k1 = 0; k1 < radix; k1++) {
      const sinkIndex = k2 + 4 * k1;
      edges.push({
        id:   `edge-col${k2}out${k1}-sink${sinkIndex}`,
        from: { node: `col-${k2}`,     port: `out${k1}` },
        to:   { node: `snk${sinkIndex}`, port: 'in' },
        delay: 600,
      });
    }
  }

  return { nodes, edges };
}
\end{lstlisting}

% -------------------------------------------------------------------
\subsection*{А.4 Короткий опис модуля ядра симуляції (src/core/scheduler.ts)}
% -------------------------------------------------------------------

Модуль \texttt{src/core/scheduler.ts} реалізує клас \texttt{DataflowRuntime} ---
ядро виконання dataflow-моделі.
Конструктор приймає об'єкт \texttt{Graph} (вузли та дуги) і функцію
\texttt{getTime(): number}, що повертає поточний логічний час симуляції.
При ініціалізації для кожної дуги графа створюється окрема FIFO-черга
типу \texttt{Map<edgeId, Token[]>}, а для кожного вузла --- запис у
\texttt{Map<nodeId, number>} (час готовності \texttt{readyAt}).

Основний метод \texttt{tick(dt: number, maxFires: number)} виконує чотири
кроки на кожному кроці симуляції:
\begin{enumerate}
  \item \textbf{Доставка токенів.} Для кожної дуги перевіряється умова
    $t_{\text{now}} - t_{\text{token}} \geq \mathtt{edge.delay}$; токени, що
    задовольняють умову, переміщуються з черги дуги до вхідного буфера
    \texttt{inbox} відповідного вузла.
  \item \textbf{Перевірка готовності вузла.} Вузол вважається готовим до
    виконання, якщо всі його вхідні порти заповнені токенами і поточний
    логічний час не менший за \texttt{readyAt[nodeId]}.
  \item \textbf{Активація вузла.} Для готового вузла викликається метод
    \texttt{execute()}, що обчислює вихідні токени залежно від типу вузла
    (\texttt{NodeKind}): \texttt{dft4} викликає \texttt{computeDFT4()},
    \texttt{twiddle} --- комплексне множення на $W_N^k$, \texttt{sink} ---
    зберігає результат у масиві виходів.
    Значення \texttt{readyAt[nodeId]} оновлюється на
    $t_{\text{now}} + \mathtt{node.latency}$.
  \item \textbf{Постановка вихідних токенів у черги.} Метод
    \texttt{emitFrom()} поміщає кожен вихідний токен у FIFO-черги дуг,
    що виходять з активованого вузла, з міткою часу $t_{\text{emit}}$.
\end{enumerate}

Клас підтримує механізм подій: \texttt{onToken} (токен доставлено на вхід
вузла), \texttt{onFire} (вузол активовано), \texttt{onOutput} (результат
з'явився на виході \texttt{sink}).
Ці події використовуються компонентом \texttt{GraphView} для анімації
руху токенів і оновлення теплової карти активності вузлів.

Повний код модуля наведено у репозиторії проекту.


\end{document}
